% (User input window)

\paragraph{Mass tower decomposition and coarse deconvolution of scale kernel: finite-order local operator, tail kernel moment and bandwidth well-determined}
The following entry implements the \(\sinh\) product structure on the Fourier side of the logistic scale kernel into an executable mass tower convolution model,
From this, a family of finite-order local deconvolution operators and their stability certificates under a bandwidth-controlled system are derived;
The computable moments and fourth-order cumulants of the tail core give the observable scale of the minimum non-Gaussian deviation, thereby promoting the first-trigger mechanism of \(q=4\) to structural necessity.

\begin{conclusion}[Infinite mass tower decomposition of scale kernel: logistic kernel is equivalent to independent Laplace cascade system]\label{con:xi-terminal-scale-kernel-mass-tower-laplace}
Let the scale unimodal kernel be
\[
\phi(s):=\frac{1}{4\cosh^{2}(s/2)}=\frac{e^{s}}{(1+e^{s})^{2}}\qquad(s\in\RR).
\]
Define the Laplace kernel for every integer \(n\ge 1\)
\[
g_n(s):=\frac{n}{2}e^{-n|s|},
\qquad
\widehat g_n(\xi)=\frac{1}{1+\xi^{2}/n^{2}}.
\]
Then there exists a family of independent random variables \((X_n)_{n\ge 1}\) (the density of \(X_n\) is \(g_n\)), such that the series
\[
Y:=\sum_{n\ge 1}X_n
\]
Converges in the sense of \(L^{2}\) and almost everywhere, and the density of their sum is strictly equal to \(\phi\).
Equivalently, let \(\phi_{\le N}:=g_1*\cdots*g_N\), then \(\phi_{\le N}\to\phi\) (\(L^2\) and weak convergence), so in the distribution sense it can be recorded as
\[
\phi=g_1*g_2*\cdots .
\]
This decomposition interprets \(\phi\) as a layer-by-layer convolution seal of an "integer mass spectrum" \(\{n\}_{n\ge 1}\): each layer \(g_n\) is the Green kernel of the one-dimensional local operator \((1-\frac{1}{n^{2}}\partial_s^{2})^{-1}\), and the product spectrum of the entire mass tower is exactly closed as \(\phi\).
\end{conclusion}
\begin{proof}[Proof sketch]
From the Fourier closed form \(\widehat\phi(\xi)=\pi\xi/\sinh(\pi\xi)\) and the Hadamard product of \(\sinh\) of the proposition \ref{prop:xi-bulk-curvature-tomography-identifiability}
\(
\sinh(\pi\xi)=\pi\xi\prod_{n\ge 1}(1+\xi^{2}/n^{2})
\)
have to
\(
\widehat\phi(\xi)=\prod_{n\ge 1}\widehat g_n(\xi)
\)。
Taking the independent random variable \(X_n\) (density is \(g_n\)), then its characteristic function is \(\widehat g_n\).
Let the partial sum \(Y_N=\sum_{n=1}^{N}X_n\), then the characteristic function of \(Y_N\) is \(\prod_{n=1}^{N}\widehat g_n\), so the density of \(Y_N\) is \(\phi_{\le N}\).
And because \(\sum_{n\ge 1}\Var(X_n)=2\sum_{n\ge 1}n^{-2}<\infty\), we get that \((Y_N)\) is Cauchy in \(L^2\), and the Kolmogorov criterion shows that it converges almost everywhere.
The characteristic function of the limit variable \(Y\) is \(\widehat\phi\), so its density is \(\phi\). Certification completed.
\end{proof}

\begin{conclusion}[Partial deconvolution is the resolution knob: the finite-order local operator reduces infinite deconvolution to tail kernel convolution]\label{con:xi-terminal-scale-kernel-partial-deconvolution-resolution}
Define finite-order local operators (understood as Fourier multipliers) for \(N\ge 1\)
\[
\mathcal D_N:=\prod_{n=1}^{N}\Bigl(1-\frac{1}{n^2}\partial_s^2\Bigr),
\qquad
\widehat{(\mathcal D_N f)}(\xi)=M_N(\xi)\,\widehat f(\xi),
\qquad
M_N(\xi):=\prod_{n=1}^{N}(1+\xi^{2}/n^{2}).
\]
And let the tail kernel \(\phi_{>N}\) be the tail sum in the conclusion \ref{con:xi-terminal-scale-kernel-mass-tower-laplace}
\(
Y_{>N}:=\sum_{n>N}X_n
\)
density (i.e. \(\widehat{\phi_{>N}}(\xi)=\prod_{n>N}(1+\xi^{2}/n^{2})^{-1}\)).
Then for any finite symbolic measure \(\nu\), if \(p=\phi*\nu\), a strict identity is established
\[
\boxed{\;
\mathcal D_N p=\phi_{>N}*\nu.\;}
\]
Therefore \(\mathcal D_N\) institutionally implements the coarse deconvolution of "removing the first \(N\) mass shell": when \(N\to\infty\) \(\phi_{>N}\Rightarrow\delta_0\), thus \(\mathcal D_N p\Rightarrow \nu\) holds in the distribution sense.
This conclusion gives an executable continuum: deconvolution replaces the "accurate or collapse" binary mechanism with a resolution stream explicitly regulated by \(N\).
\end{conclusion}
\begin{proof}[Proof sketch]
From the conclusion \ref{con:xi-terminal-scale-kernel-mass-tower-laplace}, \(\widehat\phi(\xi)=M_N(\xi)^{-1}\widehat{\phi_{>N}}(\xi)\).
Taking the Fourier transform of \(p=\phi*\nu\) we get \(\widehat p=\widehat\phi\cdot\widehat\nu\), thus
\[
\widehat{(\mathcal D_N p)}(\xi)=M_N(\xi)\widehat p(\xi)=\widehat{\phi_{>N}}(\xi)\widehat\nu(\xi)=\widehat{(\phi_{>N}*\nu)}(\xi),
\]
Get the identity.
And since \(\Var(Y_{>N})\to 0\), therefore \(\phi_{>N}\Rightarrow\delta_0\), thus \(\phi_{>N}*\nu\Rightarrow\nu\). Certification completed.
\end{proof}

\begin{conclusion}[Moments of the tail kernel are fully computable: resolution scale \texorpdfstring{$\sigma_N\sim \sqrt{2/N}$}{sigma\_N ~ sqrt(2/N)} with Wasserstein control]\label{con:xi-terminal-scale-kernel-tail-moments-w2}
Following the conclusion \ref{con:xi-terminal-scale-kernel-mass-tower-laplace} of the random variable \(X_n\), let \(Y_{>N}:=\sum_{n>N}X_n\), its density is \(\phi_{>N}\).
Then there is a certificate that can add moments
\[
\Var(Y_{>N})=\sum_{n>N}\Var(X_n)=2\sum_{n>N}\frac{1}{n^2}=: \sigma_N^2,
\qquad
\kappa_4(Y_{>N})=\sum_{n>N}\kappa_4(X_n)=12\sum_{n>N}\frac{1}{n^4}.
\]
And for any probability measure \(\nu\) with finite second moment, there is a certifiable upper bound \(W_2\)
\[
\boxed{\;
W_2\bigl(\phi_{>N}*\nu,\nu\bigr)\ \le\ \sigma_N
\ =\ \sqrt{2\sum_{n>N}n^{-2}}
\ \sim\ \sqrt{\frac{2}{N}}\quad(N\to\infty).\;}
\]
Therefore, if the resolution \(\sigma>0\) is allowed, taking \(N\asymp 2/\sigma^{2}\) can guarantee \(W_2\)-error \(\le \sigma\);
Under this regime, the resolution cost drops from exponential to power-law.
\end{conclusion}
\begin{proof}[Proof sketch]
For \(X_n\sim g_n\), we can directly calculate \(\Var(X_n)=2/n^2\), \(\kappa_4(X_n)=12/n^4\), and the cumulants are independent and additive, thus obtaining the moment certificate.
For the upper bound of \(W_2\), take the coupling: let \(U\sim \nu\), let \(V\sim \phi_{>N}\) and be independent, then \((U+V,U)\) is a coupling of \((\phi_{>N}*\nu,\nu)\), so
\(
W_2^2(\phi_{>N}*\nu,\nu)\le \EE[|V|^2]=\Var(V)=\sigma_N^2
\)。
The asymptotic \(\sum_{n>N}n^{-2}\sim 1/N\) is obtained by the integral discriminant method. Certification completed.
\end{proof}

\begin{conclusion}[Gaussianization and fourth-order head breaking of the tail kernel: non-Gaussianity is uniquely dominated by \texorpdfstring{$q=4$}{q=4} magnitude]\label{con:xi-terminal-scale-kernel-tail-gaussianization-kurtosis}
Let \(\sigma_N^2=\Var(Y_{>N})\), and define the standardized tail variable \(Z_N:=Y_{>N}/\sigma_N\).
Then there is weak convergence \(Z_N\Rightarrow \mathcal N(0,1)\), and its normalized excess kurtosis has a closed form
\[
\boxed{\;
\gamma_2(N):=\frac{\kappa_4(Y_{>N})}{\sigma_N^4}
=3\,\frac{\sum_{n>N}n^{-4}}{\bigl(\sum_{n>N}n^{-2}\bigr)^2}
\sim \frac{1}{N}\quad(N\to\infty).\;}
\]
And all odd-order cumulants are strictly zero.
The first non-Gaussian deviation of the tail core therefore inevitably occurs at the fourth order: in the observable correction of the scale flow, the minimum broken order is structurally locked to \(4\) and enters with a strength of \(N^{-1}\).
\end{conclusion}
\begin{proof}[Proof sketch]
Since \(X_n\) is symmetric, all odd-order cumulants are zero and remain zero under independent sums.
The central limit theorem for \(Z_N\) is given by the Lindeberg condition: \(\max_{n>N}\Var(X_n)/\sigma_N^2\to 0\) and \(\sum_{n>N}\Var(X_n)=\sigma_N^2\).
The kurtosis formula is obtained by substituting the cumulant certificate of the conclusion \ref{con:xi-terminal-scale-kernel-tail-moments-w2} and simplifying it;
The asymptotics are given by \(\sum_{n>N}n^{-4}\sim (3N^3)^{-1}\), \(\sum_{n>N}n^{-2}\sim N^{-1}\). Certification completed.
\end{proof}

\begin{conclusion}[Upper bound on the condition number under bandwidth constraints: Coarse deconvolution is uniformly well-determined at any fixed bandwidth]\label{con:xi-terminal-scale-kernel-bandwidth-conditioned}
Let the observation be \(\widetilde p=p+\eta\), and assume that the noise satisfies within the frequency band \(|\xi|\le\Omega\)
\[
\sup_{|\xi|\le\Omega}|\widehat\eta(\xi)|\le\varepsilon.
\]
Then for any \(N\ge 1\), the output error of the finite-order operator satisfies in the same frequency band
\[
\sup_{|\xi|\le\Omega}\bigl|\widehat{(\mathcal D_N\eta)}(\xi)\bigr|
\le
\varepsilon\cdot \sup_{|\xi|\le\Omega}M_N(\xi)
\le
\varepsilon\cdot \exp\!\Bigl(\Omega^2\sum_{n=1}^{\infty}\frac{1}{n^2}\Bigr)
=\varepsilon\cdot \exp\!\Bigl(\frac{\pi^2}{6}\Omega^2\Bigr).
\]
Therefore within any fixed bandwidth \(\Omega\), the noise amplification factor of \(\mathcal D_N\) is uniformly bounded for \(N\);
The pathological nature of coarse deconvolution does not come from "increasing order", but from "trying to achieve the \(\delta_0\) limit on infinite bandwidth".
This conclusion provides a strict basis for the computable phase diagram of scale flow: in a bandwidth-controlled regime, the resolution can be improved according to the \(N\sim\sigma^{-2}\) power law without triggering an exponential condition number explosion.
\end{conclusion}
\begin{proof}[Proof sketch]
Expressed by the Fourier multiplier of the conclusion \ref{con:xi-terminal-scale-kernel-partial-deconvolution-resolution},
\(
\widehat{(\mathcal D_N\eta)}=M_N\widehat\eta
\)
It holds on \(|\xi|\le\Omega\).
And for \(x\ge 0\) we have \(\log(1+x)\le x\), so
\[
\log M_N(\xi)=\sum_{n=1}^{N}\log\Bigl(1+\frac{\xi^2}{n^2}\Bigr)\le \xi^2\sum_{n=1}^{N}\frac{1}{n^2}
\le \Omega^2\sum_{n=1}^{\infty}\frac{1}{n^2}=\frac{\pi^2}{6}\Omega^2,
\]
This leads to the consistent upper bound shown. Certification completed.
\end{proof}

\begin{conclusion}[\texorpdfstring{$\Gamma$}{Gamma}-closed form and small frequency expansion: analytical completeness of local operators]\label{con:xi-terminal-scale-kernel-gamma-closed-form}
Note the finite order spectral multiplier
\[
M_N(\xi):=\prod_{n=1}^{N}(1+\xi^2/n^2).
\]
Then \(M_N\) has strict \(\Gamma\)-ratio closed form
\[
\boxed{\;
M_N(\xi)=\frac{\Gamma(N+1+\mathrm{i}\xi)\Gamma(N+1-\mathrm{i}\xi)}{(N!)^2\,\Gamma(1+\mathrm{i}\xi)\Gamma(1-\mathrm{i}\xi)}.\;}
\]
Thereby there is an auditable expansion under \(|\xi|=o(N)\)
\[
\log M_N(\xi)=\xi^2\sum_{n=1}^{N}\frac{1}{n^2}-\frac{\xi^4}{2}\sum_{n=1}^{N}\frac{1}{n^4}+O\!\Bigl(\xi^6\sum_{n=1}^{N}\frac{1}{n^6}\Bigr),
\]
It reveals the strict hierarchy of "the second order (\(\zeta(2)\)) determines the main amplification, and the fourth order (\(\zeta(4)\)) determines the primary correction".
This closed form enables the analytic structure of \(\mathcal D_N\) to be recalculated in finite steps, and allows a falsifiable optimal truncation of \(N\) when the bandwidth and accuracy are given.
\end{conclusion}
\begin{proof}[Proof sketch]
Set \(M_N(\xi)=\prod_{n=1}^{N}(n^2+\xi^2)/n^2=\prod_{n=1}^{N}(n+\mathrm{i}\xi)(n-\mathrm{i}\xi)/(N!)^2\), and then use the finite product identity
\(
\prod_{n=1}^{N}(n+z)=\Gamma(N+1+z)/\Gamma(1+z)
\)
That is, we get \(\Gamma\)-ratio closed form.
For small frequency expansion, use \(\log(1+u)=u-\frac{u^2}{2}+O(u^3)\) and sum \(u=\xi^2/n^2\). Certification completed.
\end{proof}

\begin{conclusion}[Two-regime recoverability of scale flows: Regime separation of power-law resolution and exponential accuracy]\label{con:xi-terminal-scale-kernel-two-regime-recoverability}
For any \(\sigma>0\), take \(N(\sigma)\) such that \(\sigma_{N(\sigma)}\le\sigma\) (where \(\sigma_N\) is the conclusion \ref{con:xi-terminal-scale-kernel-tail-moments-w2}).
Then from the conclusion \ref{con:xi-terminal-scale-kernel-partial-deconvolution-resolution} we have
\[
\mathcal D_{N(\sigma)}p=\phi_{>N(\sigma)}*\nu,
\]
A stable reconstruction with a resolution of \(\sigma\) is given; the cost is the local operator order \(2N(\sigma)\asymp\sigma^{-2}\), and the condition number is uniformly controllable under a fixed bandwidth (Conclusion \ref{con:xi-terminal-scale-kernel-bandwidth-conditioned}).
\par
On the contrary, if the precise limit of \(\sigma\downarrow 0\) is required (that is, the bandwidth-free limit of \(\phi_{>N}\Rightarrow\delta_0\)), it will inevitably enter the exponential ill-posed region: the tail core approaches \(\delta_0\) and requires simultaneous control of any high frequency, and the system must allow the bandwidth \(\Omega\to\infty\).
More precisely, the limit multiplier is obtained by \(\sinh(\pi\xi)=\pi\xi\prod_{n\ge 1}(1+\xi^2/n^2)\)
\[
M_\infty(\xi):=\lim_{N\to\infty}M_N(\xi)=\frac{\sinh(\pi\xi)}{\pi\xi},
\qquad
M_\infty(\xi)\sim \frac{1}{2\pi|\xi|}\,e^{\pi|\xi|}\quad(|\xi|\to\infty).
\]
Therefore, if "accurate" is understood as needing to approximate \(\delta_0\) on \(|\xi|\lesssim \Omega\), then the noise amplification increases exponentially with \(\Omega\);
If the system resolution \(\sigma\) is further regarded as the required visible bandwidth \(\Omega\asymp \sigma^{-1}\), then the magnitude of the amplification factor is \(\exp(\pi/\sigma)\) (only the difference is the polynomial factor), so the number of precision bits increases linearly with \(1/\sigma\).
Therefore, "power law controllable" and "exponential uncontrollable" are not contradictory, but two systems separated by scale selection (whether to accept \(\sigma\)-resolution).
\end{conclusion}
\begin{proof}[Proof sketch]
The stable reconstruction part is given by the identity of the conclusion \ref{con:xi-terminal-scale-kernel-partial-deconvolution-resolution} and the \(\sigma_N\) control of the conclusion \ref{con:xi-terminal-scale-kernel-tail-moments-w2};
Wellness in a bandwidth-controlled regime is given by the conclusion \ref{con:xi-terminal-scale-kernel-bandwidth-conditioned}.
The exact limit requires that \(\Omega\) is unbounded, so the noise amplification must diverge with \(\Omega\), giving an exponential ill-posed region. Certification completed.
\end{proof}

\begin{conclusion}[\texorpdfstring{$(q=4)$}{(q=4)}'s regularization status: the fourth-order violation determines the minimum moment order of the visible phase transition]\label{con:xi-terminal-scale-kernel-q4-minimal-visible-order}
In the mass tower model with the conclusion \ref{con:xi-terminal-scale-kernel-tail-gaussianization-kurtosis}--\ref{con:xi-terminal-scale-kernel-tail-gaussianization-kurtosis}, all odd-order cumulants of the tail core are strictly zero, and the first non-Gaussian term is given by the fourth-order cumulant \(\kappa_4\) and given by
\(\gamma_2(N)\sim N^{-1}\)
control.
Therefore, any auditable inference system based on "second-order approximation (Gaussian) + bandwidth-controlled rough deconvolution" will have its minimum systematic deviation at the fourth order;
In other words, if only finite moment orders are allowed as statistically visible interfaces, the minimum moment order capable of detecting "non-Gaussian boundaries/phase transition criticalities" must be \(4\).
This fact provides a completely structural explanation for the first line crossing phenomenon of \(q=4\): it is not an accidental value, but the minimum breaking order forced by the symmetry of the tail of the scale core and the superposition of mass towers.
\end{conclusion}
\begin{proof}[Proof sketch]
From the conclusion \ref{con:xi-terminal-scale-kernel-tail-gaussianization-kurtosis}, the odd-order cumulant of the standardized tail variable is zero, and the fourth-order cumulant gives the first non-Gaussian deviation and decays with \(N^{-1}\) magnitude.
In a system based on the second-order Gaussian approximation, the systematic deviation of any finite-order moment statistic is determined by the lowest non-zero cumulant order, so the minimum visible moment order is \(4\). Certification completed.
\end{proof}

% (User input window)
%
% (Source tag) /killo-golden

\paragraph{Tree/ring layering of determinant volume entropy: Chow--Liu projection, conditional mutual information and fourth-order development}
Under the caliber of carrying the "volume potential" with the Gram/correlation matrix,\(-\log\det\) naturally gives a many-body volume entropy potential;
This paragraph compresses the auditable lower bound of this potential into a Chow--Liu tree projection and interprets the remaining gap as an irreversible volume loss carried by the ring structure.
In the one-dimensional sorting and monotonic distance kernel system, the optimal tree selection is rigidly degenerated into adjacent chains, thereby externalizing the "volume entropy carried by the tree" into an adjacent interval ledger;
In the weak coherence region, the first systematic contribution of ring entropy is controlled by conditional mutual information and develops from the fourth order, which is consistent with what appears repeatedly in this article. \(q=4\) The first trigger structure is of the same type.

\begin{conclusion}[\texorpdfstring{$(\det)$}{(det)}-Chow--Liu tree lower bound of entropy potential: the determinant volume is controlled by the maximum spanning tree]\label{con:xi-det-entropy-chow-liu-tree-lower-bound}
make \(R\in\RR^{\kappa\times\kappa}\) is the correlation matrix (\(R\succ0\)、\(\mathrm{diag}(R)=\mathbf 1\)), define the determinant volume entropy potential
\[
S_{\det}(R):=-\log\det R.
\]
For each undirected edge \((i,j)\) Define the two-body mutual information weight
\[
w_{ij}:=-\log(1-R_{ij}^2).
\]
Then for any spanning tree \(T\) have
\[
S_{\det}(R)\ \ge\ \sum_{(i,j)\in T} w_{ij},
\qquad\text{thereby}\qquad
S_{\det}(R)\ \ge\ \max_{T}\sum_{(i,j)\in T} w_{ij}.
\]
and for the given \(T\), the equal sign holds if and only if \(R^{-1}\), the non-diagonal non-zero support is contained in \(T\) (equivalent to \(\mathcal N(0,R)\) in the picture \(T\) is a tree-shaped Gaussian Markov field).
\end{conclusion}
\begin{proof}[Points of proof]
Pick \(X\sim\mathcal N(0,R)\)。
in fixed tree \(T\), Chow--Liu tree projection \(\mathcal N_T\) (see conclusion \ref{con:xi-det-entropy-loop-entropy-kl-gap} settings) satisfy
\(
2D_{\mathrm{KL}}(\mathcal N(0,R)\,\|\,\mathcal N_T)
=S_{\det}(R)-\sum_{(i,j)\in T}w_{ij}\ge 0
\)，
Thus the lower realm is obtained.
The equal sign is equivalent to \(D_{\mathrm{KL}}=0\),Right now \(\mathcal N(0,R)=\mathcal N_T\)；
For the Gaussian family this is equivalent to the accuracy matrix only in \(T\) is allowed to have non-zero terms on its edges. Certification completed.
\end{proof}

\begin{conclusion}[Ring entropy is the tree approximation \texorpdfstring{$D_{\mathrm{KL}}$}{DKL} Gap: Exact tree decomposition of volumetric losses]\label{con:xi-det-entropy-loop-entropy-kl-gap}
Follow the conclusion \ref{con:xi-det-entropy-chow-liu-tree-lower-bound} settings.
For a given spanning tree \(T\) defines tree ring entropy
\[
\mathcal L_T(R):=S_{\det}(R)-\sum_{(i,j)\in T}w_{ij}
=-\log\det R+\sum_{(i,j)\in T}\log(1-R_{ij}^2).
\]
but \(\mathcal L_T(R)\ge 0\) is always established.
like \(X\sim\mathcal N(0,R)\), and order \(\mathcal N_T\) think \(T\) is the Chow--Liu tree Gaussian projection of the graph (its one-dimensional edge and tree edge binary edge are \(\mathcal N(0,R)\) are consistent), then the strict identity is established
\[
\mathcal L_T(R)=2D_{\mathrm{KL}}\!\bigl(\mathcal N(0,R)\,\|\,\mathcal N_T\bigr).
\]
therefore
\(
\arg\min_T\mathcal L_T(R)=\arg\max_T\sum_{(i,j)\in T}w_{ij}
\)
Gives the optimal tree approximation.
\end{conclusion}
\begin{proof}[Points of proof]
make \(R_T\) for \(\mathcal N_T\) correlation matrix.
tree projection \(\mathcal N_T\) is uniquely determined by the tree edge correlation and satisfies
\[
\det R_T=\prod_{(i,j)\in T}(1-R_{ij}^2)
\]
(Edge-wise conditional variance product formula for tree-shaped Gaussians).
On the other hand, the closed form of KL for the family of zero-mean Gaussians gives
\[
2D_{\mathrm{KL}}(\mathcal N(0,R)\,\|\,\mathcal N(0,R_T))
=\tr(R_T^{-1}R)-\kappa+\log\frac{\det R_T}{\det R}.
\]
because \(R_T^{-1}\), the off-diagonal support is contained in \(T\),and \(R\) and \(R_T\) has the same value at the diagonal and on the edge of the tree, so
\(
\tr(R_T^{-1}R)=\tr(R_T^{-1}R_T)=\kappa
\)，
thereby
\(
2D_{\mathrm{KL}}=\log\frac{\det R_T}{\det R}
\)。
Substituting into the determinant factor, we get the identity shown. Certification completed.
\end{proof}

\begin{conclusion}[Tree selection rigidity of monotonic distance kernel: the optimal tree of one-dimensional sorting points must be adjacent chains]\label{con:xi-det-entropy-monotone-kernel-chain-mst}
Suppose there exists a strictly monotonic decreasing function \(A:(0,\infty)\to(0,1)\) make
\[
R_{ij}=A(|s_i-s_j|),
\qquad
s_1<\cdots<s_\kappa.
\]
Then the edge right \(w_{ij}=-\log(1-A(|s_i-s_j|)^2)\) also strictly decreases with distance.
Therefore, the maximum spanning tree is uniquely a path graph (adjacent chain)
\[
T_{\mathrm{chain}}:=\{(1,2),(2,3),\dots,(\kappa-1,\kappa)\},
\]
So as to get the lower bound of the strongest tree
\[
S_{\det}(R)\ \ge\ \sum_{j=1}^{\kappa-1}-\log\!\bigl(1-A(\Delta_j)^2\bigr),
\qquad
\Delta_j:=s_{j+1}-s_j.
\]
\end{conclusion}
\begin{proof}[Points of proof]
For any cut \(\{1,\dots,j\}\,\cup\,\{j+1,\dots,\kappa\}\), the smallest edge across the cut distance is uniquely \((j,j+1)\), and its weight is the largest.
According to the cut property of the maximum spanning tree, each \((j,j+1)\) must belong to the maximum spanning tree; merge all \(j\) to obtain uniqueness and chain structure. Certification completed.
\end{proof}

\begin{conclusion}[Closed form of three-point local ring entropy: the chain gap is equal to the conditional mutual information and is a partially correlated logarithmic potential]\label{con:xi-det-entropy-three-point-loop-entropy-partial-corr}
Take any \(a<b<c\),remember \(\rho_{ab}=R_{ab}\)、\(\rho_{bc}=R_{bc}\)、\(\rho_{ac}=R_{ac}\), and order \(R_{abc}\) is the corresponding three-element correlation submatrix.
Relative to chain tree \((a,b)\)--\((b,c)\) Define three-point local ring entropy
\[
\mathcal L^{(3)}_{a,b,c}
:=-\log\det R_{abc}+\log(1-\rho_{ab}^2)+\log(1-\rho_{bc}^2).
\]
Then there is the identity
\[
\mathcal L^{(3)}_{a,b,c}
=-\log\bigl(1-\rho_{ac\mid b}^2\bigr),
\qquad
\rho_{ac\mid b}:=\frac{\rho_{ac}-\rho_{ab}\rho_{bc}}{\sqrt{(1-\rho_{ab}^2)(1-\rho_{bc}^2)}}.
\]
And if \(X\sim\mathcal N(0,R)\),but
\[
\mathcal L^{(3)}_{a,b,c}=2I(X_a;X_c\mid X_b).
\]
\end{conclusion}
\begin{proof}[Points of proof]
Standard identities based on determinants and partial correlations of ternary correlation matrices
\(
1-\rho_{ac\mid b}^2=\det R_{abc}/((1-\rho_{ab}^2)(1-\rho_{bc}^2))
\)
Get the first formula.
For the Gaussian family, the conditional mutual information satisfies
\(
2I(X_a;X_c\mid X_b)=-\log(1-\rho_{ac\mid b}^2)
\)，
Thus we get the second formula. Certification completed.
\end{proof}

\begin{conclusion}[The fourth-order first-trigger law: the ring entropy of the long-distance attenuation weak coherence region begins to develop from the fourth order]\label{con:xi-det-entropy-q4-trigger-weak-coherence}
in conclusion \ref{con:xi-det-entropy-three-point-loop-entropy-partial-corr} Under the three-point settings, if
\(
|\rho_{ab}|,|\rho_{bc}|,|\rho_{ac}|\le\varepsilon\ll1
\)，
then there is a consistent expansion
\[
\rho_{ac\mid b}=(\rho_{ac}-\rho_{ab}\rho_{bc})+O(\varepsilon^3),
\qquad
\mathcal L^{(3)}_{a,b,c}=\rho_{ac\mid b}^2+O(\varepsilon^4).
\]
In particular, when the long-distance attenuation type relationship is satisfied \(\rho_{ac}=\rho_{ab}\rho_{bc}+O(\varepsilon^3)\) hour,\(\rho_{ac\mid b}=O(\varepsilon^2)\) and
\[
\mathcal L^{(3)}_{a,b,c}=O(\varepsilon^4),
\]
That is, the local ring entropy has "minimum fourth-order visibility": two-body correlation \(O(\varepsilon)\) is not sufficient to trigger a chain gap under this regime, the first systematic magnitude of the gap must be fourth order.
\end{conclusion}
\begin{proof}[Points of proof]
Denominator \(\sqrt{(1-\rho_{ab}^2)(1-\rho_{bc}^2)}=1+O(\varepsilon^2)\), so part of the correlation can be obtained by linearizing the molecule.
Also by \(-\log(1-x)=x+O(x^2)\)（\(x\to0\)) and take \(x=\rho_{ac\mid b}^2=O(\varepsilon^2)\), expanded as shown. Certification completed.
\end{proof}

\begin{conclusion}[Explicit fourth-order scaling under exponential tail scaling kernel:\texorpdfstring{$\mathcal L^{(3)}$}{L(3)} auditable asymptotic]\label{con:xi-det-entropy-exp-tail-kernel-audit-asymptotics}
Take a specific monotonic distance kernel
\[
A(\Delta):=\frac{\Delta}{2\sinh(\Delta/2)}\qquad(\Delta\ge0),
\qquad
A(0):=1,
\]
And order \(\Delta_1=s_b-s_a\)、\(\Delta_2=s_c-s_b\)。
when \(\Delta_1,\Delta_2\to\infty\) sometimes
\[
\rho_{ab}\sim \Delta_1e^{-\Delta_1/2},\quad
\rho_{bc}\sim \Delta_2e^{-\Delta_2/2},\quad
\rho_{ac}\sim (\Delta_1+\Delta_2)e^{-(\Delta_1+\Delta_2)/2}.
\]
thereby
\[
\rho_{ac}-\rho_{ab}\rho_{bc}
=e^{-(\Delta_1+\Delta_2)/2}\Bigl((\Delta_1+\Delta_2)-\Delta_1\Delta_2+o(\Delta_1\Delta_2)\Bigr),
\]
and obtain auditable fourth-order dominant terms in the weak coherence region
\[
\mathcal L^{(3)}_{a,b,c}
=\bigl(\rho_{ac}-\rho_{ab}\rho_{bc}\bigr)^2\,(1+o(1))
\asymp (\Delta_1\Delta_2)^2\,e^{-(\Delta_1+\Delta_2)}.
\]
\end{conclusion}
\begin{proof}[Points of proof]
Depend on \(A(\Delta)\sim \Delta e^{-\Delta/2}\)（\(\Delta\to\infty\)) Three correlations are asymptotically obtained.
Reuse conclusion \ref{con:xi-det-entropy-q4-trigger-weak-coherence} Will \(\mathcal L^{(3)}\) linearized to \(\rho_{ac\mid b}^2\), and note the denominator \(1+o(1)\), the dominant term and scale can be obtained. Certification completed.
\end{proof}

\begin{conclusion}[Chain decomposition of chain tree ring entropy: the global gap is equal to the sum of far past conditional mutual information]\label{con:xi-det-entropy-chain-loop-entropy-chain-rule}
make \(X=(X_1,\dots,X_\kappa)\sim\mathcal N(0,R)\), and take the adjacent chain tree
\(
T_{\mathrm{chain}}=\{(1,2),\dots,(\kappa-1,\kappa)\}
\)。
Define global chain entropy
\[
\mathcal L_{\mathrm{chain}}(R):=\mathcal L_{T_{\mathrm{chain}}}(R)
=-\log\det R+\sum_{j=1}^{\kappa-1}\log(1-R_{j,j+1}^2).
\]
Then there is an exact chain decomposition
\[
\mathcal L_{\mathrm{chain}}(R)
=2\sum_{j=3}^{\kappa} I\!\bigl(X_j;X_{1:j-2}\mid X_{j-1}\bigr),
\]
thereby
\[
\mathcal L_{\mathrm{chain}}(R)\ \ge\ 2\sum_{j=3}^{\kappa} I\!\bigl(X_j;X_{j-2}\mid X_{j-1}\bigr)
=\sum_{j=3}^{\kappa}\mathcal L^{(3)}_{j-2,j-1,j}.
\]
\end{conclusion}
\begin{proof}[Points of proof]
By conclusion \ref{con:xi-det-entropy-loop-entropy-kl-gap}，\(\mathcal L_{\mathrm{chain}}/2\) is equal to \(\mathcal N(0,R)\) KL gap relative to the chain Chow--Liu projection;
For a chain graph the gap is equal to the sum of a list of conditional mutual information (chain rule).
Lower bound by term-wise nonnegative decomposition of conditional mutual information
\(
I(X_j;X_{1:j-2}\mid X_{j-1})=\sum_{i=1}^{j-2}I(X_j;X_i\mid X_{i+1:j-1})\ge I(X_j;X_{j-2}\mid X_{j-1})
\)
given. Certification completed.
\end{proof}

\begin{conclusion}[of weak coherence chain \(q=4\) Volume law: The dominant term of ring entropy is the sum of squares of second-order differential deviations]\label{con:xi-det-entropy-weak-coherence-chain-q4-law}
in conclusion \ref{con:xi-det-entropy-monotone-kernel-chain-mst} Under the setting of adjacent sorting, note \(\rho_{j,j+1}=R_{j,j+1}\) and define the second-order differential deviation
\[
\eta_j:=R_{j-2,j}-R_{j-2,j-1}R_{j-1,j}\qquad(3\le j\le\kappa).
\]
If the whole is in a weak coherence region (for example, taking all adjacent intervals for the exponential tail core \(\Delta_j\) is large enough that all \(|R_{ij}|\) is uniformly small enough), then there is a consistent approximation
\[
\mathcal L_{\mathrm{chain}}(R)
=\sum_{j=3}^{\kappa}\frac{\eta_j^2}{(1-\rho_{j-2,j-1}^2)(1-\rho_{j-1,j}^2)}
\ +\ o\!\Bigl(\sum_{j=3}^{\kappa}\eta_j^2\Bigr).
\]
The first systematic contribution of the chain-link entropy is therefore controlled by the sum of squares of the "second-order differential deviations" and exhibits fourth-order dominant magnitude in the long-distance decay regime.
\end{conclusion}
\begin{proof}[Points of proof]
By conclusion \ref{con:xi-det-entropy-chain-loop-entropy-chain-rule}，\(\mathcal L_{\mathrm{chain}}\) is given by a list of conditional mutual information.
In the weak coherence region, the far past pair \(X_j\) conditional influence on higher order can be ignored, and the dominant term is given by \(I(X_j;X_{j-2}\mid X_{j-1})\) given;
and by conclusion \ref{con:xi-det-entropy-three-point-loop-entropy-partial-corr} have to
\(
\mathcal L^{(3)}_{j-2,j-1,j}=-\log(1-\rho_{j-2,j\mid j-1}^2)
\)
and
\(
\rho_{j-2,j\mid j-1}^2=\eta_j^2/((1-\rho_{j-2,j-1}^2)(1-\rho_{j-1,j}^2))
\)。
reuse \(-\log(1-x)=x+o(x)\) can be linearized. Certification completed.
\end{proof}

\begin{conclusion}[The ratio of Markov volume to real volume: the link entropy is the exact logarithm of the volume deficit]\label{con:xi-det-entropy-markov-volume-ratio}
Define chain Markov correlation matrix \(R^{\mathrm{MC}}\) to satisfy
\[
R^{\mathrm{MC}}_{jj}=1,\qquad
R^{\mathrm{MC}}_{j,j+1}=R_{j,j+1},
\qquad
(R^{\mathrm{MC}})^{-1}\ \text{for three diagonals},
\]
The unique correlation matrix of (that is, with \(T_{\mathrm{chain}}\) compatible tree Gaussian projection).
Then the closed form of determinant is established
\[
\det R^{\mathrm{MC}}=\prod_{j=1}^{\kappa-1}\bigl(1-R_{j,j+1}^2\bigr),
\qquad
\mathcal L_{\mathrm{chain}}(R)= -\log\frac{\det R}{\det R^{\mathrm{MC}}}.
\]
The link entropy is therefore the exact logarithmic loss of the "real volume" relative to the "Markov volume", providing a scale for volume comparison without the need for external fitting.
\end{conclusion}
\begin{proof}[Points of proof]
The chain tree projection corresponds to a column of stepwise regression representation
\(
X_{j+1}=R_{j,j+1}X_j+\sqrt{1-R_{j,j+1}^2}\,\varepsilon_{j+1}
\)
（\(\varepsilon_{j+1}\) independent standard Gaussian), so that the covariance determinant is the conditional variance product, giving \(\det R^{\mathrm{MC}}\) factorization;
And then the conclusion \ref{con:xi-det-entropy-loop-entropy-kl-gap} Take the chain tree \(T=T_{\mathrm{chain}}\) to obtain the volume ratio identity. Certification completed.
\end{proof}

\begin{conclusion}[Basic ring conditional mutual information certificate for arbitrary trees: the cost of non-tree edges is governed by global ring entropy]\label{con:xi-det-entropy-tree-basecycle-cmi-certificate}
for any spanning tree \(T\) and any non-tree edge \(e=(u,v)\notin T\),make \(P_T(u,v)\) for tree connections \(u\) and \(v\) of the unique path vertex set, and remember the separation set
\(
S_e:=P_T(u,v)\setminus\{u,v\}
\)。
Define the base ring conditional mutual information potential
\[
\mathcal I_e:=2I\!\bigl(X_u;X_v\mid X_{S_e}\bigr)
=-\log\bigl(1-\rho_{uv\mid S_e}^2\bigr),
\]
in \(\rho_{uv\mid S_e}\) is the pair under the Gaussian family \(S_e\) conditional partial correlation coefficient.
Then there is always edge-by-edge dominance over the Gaussian family
\[
\mathcal I_e\ \le\ \mathcal L_T(R)\qquad\text{to all }e\notin T.
\]
Therefore, the conditional correlation strength of any non-tree edge is controlled by the same global ring entropy; under the weak coherence regime (the non-tree edge correlation is dominated by the path correlation product), each \(\mathcal I_e\) is given by the fourth order magnitude, so that \(\mathcal L_T\) is structurally locked under this system to \(4\)。
\end{conclusion}
\begin{proof}[Points of proof]
tree projection \(\mathcal N_T\) in the picture \(T\) satisfies separability: for any \(e=(u,v)\notin T\), there must be \(X_u\perp X_v\mid X_{S_e}\) At \(\mathcal N_T\) is established under.
set up \(\mathcal Q_e\) are all distribution families that satisfy the independence condition, then
\(
\mathcal N_T\in\mathcal Q_e
\)。
on the other hand,\(D_{\mathrm{KL}}(\mathcal N(0,R)\,\|\,\mathcal Q_e)\) is exactly \(I(X_u;X_v\mid X_{S_e})\)；
therefore
\(
D_{\mathrm{KL}}(\mathcal N(0,R)\,\|\,\mathcal N_T)\ge I(X_u;X_v\mid X_{S_e})
\)。
multiply by \(2\) and use the conclusion \ref{con:xi-det-entropy-loop-entropy-kl-gap} The identity of is governed as shown. Certification completed.
\end{proof}

\begin{conclusion}[Per-step ring entropy scaling for the isometric weakly coherent phase:\texorpdfstring{$\Delta\to\infty$}{Delta->inf} hour \texorpdfstring{$\mathcal L_{\mathrm{chain}}/\kappa$}{L_chain/kappa} explicit attenuation]\label{con:xi-det-entropy-equispaced-loop-entropy-scaling}
Get equidistant depth \(s_j=s_0+(j-1)\Delta\) and order \(\Delta\to\infty\)。
remember
\(
a:=A(\Delta)\sim \Delta e^{-\Delta/2}
\)、
\(
b:=A(2\Delta)\sim 2\Delta e^{-\Delta}
\)，
Then the second-order difference deviation of each adjacent three points satisfies
\[
\eta=b-a^2\sim-\Delta^2e^{-\Delta}.
\]
Therefore, the strength of each step of the chain link entropy satisfies
\[
\frac{1}{\kappa}\mathcal L_{\mathrm{chain}}(R)
\sim \eta^2
\asymp \Delta^{4}e^{-2\Delta},
\qquad (\Delta\to\infty,\ \kappa\ \text{big}).
\]
This formula externalizes the non-Markovian property under weak coherence into a fully explicit exponential law: the exponential rate \(2\Delta\), and use polynomial \(\Delta^4\) as a pre-factor correction.
\end{conclusion}
\begin{proof}[Points of proof]
By conclusion \ref{con:xi-det-entropy-exp-tail-kernel-audit-asymptotics} The kernel asymptotically obtains \(a\)、\(b\), and calculate \(\eta=b-a^2\)。
And then the conclusion \ref{con:xi-det-entropy-weak-coherence-chain-q4-law}, weak coherence region \(\mathcal L_{\mathrm{chain}}\) the dominant term is \(\sum \eta_j^2\); in the equidistant case, all terms are of the same order, divided by \(\kappa\) to get the scale for each step. Certification completed.
\end{proof}

\begin{conclusion}[Tree/ring layering of volume-dimension-entropy: boundary dimensions are given by tree terms, fractal roughness is given by ring terms]\label{con:xi-det-entropy-tree-loop-stratification}
For any correlation matrix \(R\) definition
\[
S_{\mathrm{tree}}(R):=\max_{T}\sum_{(i,j)\in T}-\log(1-R_{ij}^2),
\qquad
S_{\mathrm{loop}}(R):=S_{\det}(R)-S_{\mathrm{tree}}(R)\ (\ge0).
\]
but \(S_{\mathrm{tree}}\) gives "the strongest volumetric entropy that can be carried by a tree structure", while \(S_{\mathrm{loop}}\) gives "the additional volumetric entropy that must be carried by the ring structure".
In the case of one-dimensional sorting with monotonic distance kernel,\(S_{\mathrm{tree}}\) from the conclusion \ref{con:xi-det-entropy-monotone-kernel-chain-mst} Rigidity degenerates into adjacent interval ledgers and gives the dominant volume entropy contribution on the boundary side;
and \(S_{\mathrm{loop}}\) in the weak coherence region, it is concluded that \ref{con:xi-det-entropy-q4-trigger-weak-coherence}--\ref{con:xi-det-entropy-weak-coherence-chain-q4-law} It develops from the fourth order and can be regarded as the minimum visible volume index describing "fractal roughness/phase change criticality".
\end{conclusion}
\begin{proof}[Points of proof]
nonnegativity by conclusion \ref{con:xi-det-entropy-chow-liu-tree-lower-bound} Launch immediately.
Chain rigidity under one-dimensional ordering is given by the conclusion \ref{con:xi-det-entropy-monotone-kernel-chain-mst} given;
The fourth-order development of the weak coherence region is based on the conclusion \ref{con:xi-det-entropy-q4-trigger-weak-coherence} and conclusion \ref{con:xi-det-entropy-weak-coherence-chain-q4-law} given.
Therefore, tree items and ring items respectively correspond to two types of structural bearings. Certification completed.
\end{proof}

% (User input window)
%
% (Source tag) /killo-golden
% Timestamp: 2026-02-18 14:02:43 (Asia/Singapore)
%
% The content is rewritten using the existing notation and auditable interface of this article, without introducing manual numbering.

\paragraph{Locatable readouts and cross-aperture invariants with finite falsification: Toeplitz pencils, negative squared exponents, endpoint budgets, and arithmetic guardrails}
This paragraph gives a set of itemized reductions consistent with the auditable chain of the main text/appendix: finite-level Toeplitz failure not only provides a binary falsifier, but can also be promoted to candidate point generation, endpoint depth readout and closure semantics of exponential invariants; and in warped graphs \(\zeta\) and \(p\) - Enumerable obstacle classifications and zero-cost audit guardrails are given in the power congruence caliber.

\begin{conclusion}[Constructive readout of finite falsifiers: Toeplitz failure can be algorithmized into offline zero candidate output]\label{con:xi-finite-falsifier-offline-fiber-readout}
Follow the horizon Carath\'eodory coefficient \((c_n)\) and the Toeplitz master-submatrix \(\mathbf{T}_N=(c_{j-k})_{0\le j,k\le N}\) notation (theorem \ref{thm:app-horizon-toeplitz-lmi}）。
If there is a \(N\) make \(\mathbf{T}_N\nsucceq 0\), then the first-order shifted Toeplitz matrix can be
\[
\mathbf{T}_N^{(1)}:=(c_{j-k+1})_{0\le j,k\le N}
\]
Form Toeplitz Pencil \((\mathbf{T}_N,\mathbf{T}_N^{(1)})\), and read out the set of finite resonance candidate points within the unit disk from its generalized spectrum
\[
\Lambda_N:=\mathrm{Spec}(\mathbf{T}_N,\mathbf{T}_N^{(1)})\cap\mathbb{D}
=\{\lambda^{(N)}_j\}_{j=1}^{\kappa_N}\subset\mathbb{D},
\qquad \kappa_N:=\nu_-(\mathbf{T}_N)\ge 1.
\]
for each \(\lambda^{(N)}_j\) takes the Cayley inverse mapping \(t^{(N)}_j:=\mathcal{C}^{-1}(\lambda^{(N)}_j)\in\mathbb{H}\), and write
\(
t^{(N)}_j=\gamma^{(N)}_j+\mathrm{i}\delta^{(N)}_j
\)
（\(\gamma^{(N)}_j\in\RR\)、\(\delta^{(N)}_j>0\)）。
Then candidate offline zero-point fibers can be explicitly generated as
\[
\boxed{\;
\rho^{(N)}_j:=\frac12+\mathrm{i}\gamma^{(N)}_j-\delta^{(N)}_j\qquad(1\le j\le \kappa_N).\;}
\]
In the case of "finite type and boundary without poles" and the negative inertia index is stable on high-order truncation, the candidate set is \(N\to\infty\) in the limit converges to the set of real intra-disk poles (equivalent to offline zero-point fibers), thereby improving the finite-level Toeplitz falsification from a binary failure witness to a localizable counterexample geometric readout.
\end{conclusion}
\begin{proof}[Points of proof]
candidate point set \(\Lambda_N\) is read out by the theorem \ref{thm:xi-resonance-pencil-reconstruction} Given; fiber structure and finite type convergence aperture are given by the theorem \ref{thm:xi-finite-witness-to-offline-fiber} given.
Negative inertial stability and finite type equivalent interface see theorem \ref{thm:xi-toeplitz-inertia-stabilizes-to-kappa} and inference \ref{cor:xi-inertia-stable-counts-disk-zeros}。
\end{proof}

\begin{conclusion}[Stability law of Toeplitz's negative square exponent: negative inertia is ultimately equal to Blaschke's defect dimension]\label{con:xi-toeplitz-negative-square-stability-law}
In the finite defect setting (\(F_{\zeta}\) is a bounded type, the boundary has no poles and its Blaschke factor is a finite product), denoted
\(\kappa:=\deg(B_{F_{\zeta}})=\dim K_{B_{F_{\zeta}}}\)。
Then the disk core \(K_{\mathscr{C}_{\zeta}}\) of Pontryagin negative square is exactly \(\kappa\), and the negative inertia index of the Toeplitz master-submatrix family satisfies
\[
\nu_-(\mathbf{T}_N)\le \kappa\qquad(\forall\,N),
\qquad
\nu_-(\mathbf{T}_N)=\kappa\qquad(N\ \text{sufficiently large}).
\]
Therefore RH is equivalent to \(\kappa=0\); Any deviation at a finite level must appear as an ultimately stable finite-dimensional negative direction, thereby compressing the complexity of Toeplitz failure into a finite integer invariant.
\end{conclusion}
\begin{proof}[Points of proof]
This is a "finite pole" \(\Longleftrightarrow\) standard alignment of "finite negative square": see theorem \ref{thm:xi-toeplitz-negative-squares-equals-blaschke-defect} and theorem \ref{thm:xi-toeplitz-inertia-stabilizes-to-kappa}(and with \S\ref{subsubsec:xi-negative-square-index-pontryagin-realization} The caliber of the negative square exponent is consistent).
\end{proof}

\begin{conclusion}[Endpoint visibility lower bound of the horizon certificate order: the minimum Toeplitz failure order obeys the universal scaling law]\label{con:xi-endpoint-visibility-lower-bound-toeplitz-order}
If there is a zero point away from the critical line \(\rho=\beta+\mathrm{i}\gamma\)（\(\beta<\tfrac12\)、\(\delta=\tfrac12-\beta>0\)), remember its disc image \(w_\rho=\mathcal{C}(\gamma+\mathrm{i}\delta)\in\mathbb{D}\), and define the endpoint boundary layer depth
\[
\boxed{\ h(\rho):=1-|w_\rho|^2=\frac{4\delta}{\gamma^2+(1+\delta)^2}\ }.
\]
Then any scheme that attempts to detect the defect at the finite-order Toeplitz--PSD level has a minimum certificate level of \(N_{\min}\) must meet the quantitative threshold
\[
N_{\min}\ \gtrsim\ \frac{1}{h(\rho)}\,\log\!\bigl(1+\delta\,h(\rho)^{-1}\bigr),
\qquad\text{thereby particularly}\qquad
N_{\min}\ \gtrsim\ h(\rho)^{-1}.
\]
exist \(|\gamma|\to\infty\) and \(\delta\) is bounded,
\(
h(\rho)^{-1}\asymp \gamma^2/\delta
\)
The universal scaling law is given: the secondary radial compression of the endpoint will force the off-slice defect into the endpoint boundary layer, and lock its visibility cost at \(h(\rho)^{-1}\), cannot be circumvented by any finite depth certificate.
\end{conclusion}
\begin{proof}[Points of proof]
See the proposition for the lower bound of certificate order \ref{con:xi-toeplitz-length-depth-duality} and appendix theorem \ref{thm:app-horizon-toeplitz-detection-threshold}；
See the theorem for deep closed-form identities \ref{thm:xi-offcritical-quadratic-radial-compression}, and is given asymptotically by its high height \(h(\rho)^{-1}\asymp \gamma^2/\delta\)。
\end{proof}

\begin{conclusion}[Adams coarse-grained falsification backtesting: there is a product lower bound between the folding scale and the certificate order]\label{con:xi-adams-coarsening-backprop-budget-product}
set up \(\mathscr{C}(w)=c_0+2\sum_{n\ge 1}c_n w^n\) for Carath\'eodory The moment sequence caliber of the horizon function (and the convention \(c_{-n}=\overline{c_n}\)）。
for integers \(m\ge 1\) Define push-forward coarse-grained moments
\(
c^{(m)}_n:=c_{mn}
\)
with the derived Toeplitz hierarchy
\[
\mathbf{T}^{(m)}_N:=\bigl(c_{m(j-k)}\bigr)_{0\le j,k\le N}.
\]
Then for any \(m,N\ge 1\) has an implication relationship
\[
\boxed{\ \mathbf{T}_{mN}\succeq 0\ \Longrightarrow\ \mathbf{T}^{(m)}_{N}\succeq 0\ },
\qquad\text{Its converse proposition is}\qquad
\mathbf{T}^{(m)}_{N}\nsucceq 0\ \Longrightarrow\ \mathbf{T}_{mN}\nsucceq 0.
\]
Therefore, any falsification at the derivation level can be strictly deduced back to the falsification at the original level, so that "the deeper level dominates all its coarse-grained derivation."
Under the endpoint boundary layer mechanism, if the defect depth is \(h(\rho)\), then the detection threshold is \(mN\gtrsim h(\rho)^{-1}\) Obtain an auditable time-certificate budget product lower bound: folding scale \(m\) and horizon order \(N\) of the mixed strategy must pay \(h(\rho)^{-1}\) dimensional resource costs.
\end{conclusion}
\begin{proof}[Points of proof]
Derived level \(\mathbf{T}^{(m)}_N\succeq 0\)'s pushforward explanation is the same as \(\mathbf{T}_{mN}\Rightarrow \mathbf{T}^{(m)}_{N}\), the structural implication of \eqref{eq:xi-adams-toeplitz-psd}；
The lower bound of the product is given by the proposition \ref{con:xi-toeplitz-length-depth-duality} right \(mN\) the original level is applicable and combined \(h(\rho)\) is obtained by compressing the endpoints in closed form.
\end{proof}

\begin{conclusion}[Unification of operator value horizon positivity and strictification: complete positivity failure is equivalent to finite-dimensional falsifier residue]\label{con:xi-operator-valued-horizon-cp-falsifier}
Convert the scalar Carath\'eodory positivity is improved to take values ​​in finite dimensions \(C^\ast\)-algebra \(\mathcal{A}\) operator value Carath\'eodory class, decomposed by Peter--Weyl
\(
\mathcal{A}\cong\bigoplus_{\pi}M_{d_\pi}(\CC)
\)
Complete positivity can be strictly aligned as a Toeplitz--PSD condition on a finite number of representation blocks: operator-valued functions \(F:\mathbb{D}\to\mathcal{A}\) satisfy \(\Re F\succeq 0\) if and only if its block Toeplitz truncates \(\mathbf{T}^{\mathcal{A}}_N(F)\succeq 0\) for all \(N\) is established.
In the context of strictification, the maximum visibility quotient gives the smallest encapsulation of a completely positive visibility implementation; true resonance or non-strictifyable residuals are strictified into fully positive finite-dimensional falsifiers, thus compressing falsifiable points in the multi-channel case into PSD failures on finite representation components.
\end{conclusion}
\begin{proof}[Points of proof]
The equivalence of complete positivity and block Toeplitz--PSD is given by the theorem \ref{thm:xi-operator-herglotz-toeplitz-completeness} given;
See the theorem for channel-by-channel alignment under decomposition \ref{thm:xi-fourier-pontryagin-block-toeplitz}；
See the theorem for an explanation of the completely positive implementation of strictification \ref{thm:xi-gbc-maximal-cp-visible-realization}。
\end{proof}

\paragraph{Modular barriers, cyclotomic factors and relative spectral gaps: finite computability of resonant role groups}

\begin{conclusion}[Enumerability of modular obstacles: non-trivial twisted channels up to radius if and only if cyclotomic factors occur]\label{con:xi-mod-obstruction-cyclotomic-factor}
to ihara/kinetics \(\zeta\) of the twisted channel
\[
Z_K(u,t;\rho):=\det(I-tB_\rho(u))^{-1},
\]
There exists a nontrivial modular obstacle if and only if some \(\rho\neq \mathbf{1}\) satisfy \(\rho(B_\rho(u))=\lambda_1(u)\) (same radius as ordinary channel);
This condition is equivalent to \(\det(I-tB_\rho(u))\) contains the cyclotomic factor (that is, the unit root factor on the unit circle).
Therefore, the source of obstacles can be reduced to the cyclotomic factor classification problem of finite-dimensional integral coefficient polynomials; this discrimination is in the general distortion graph \(\zeta\) caliber is also true.
\end{conclusion}
\begin{proof}[Points of proof]
See conclusion \ref{con:xi-terminal-mod-obstruction-cyclotomic} and the classification theorem it cites \ref{thm:conclusion75-ihara-cyclotomic-factor-classification}。
\end{proof}

\begin{conclusion}[Relative spectral gap after barrier stripping: resonance only occupies a finite-dimensional identifiable subspace]\label{con:xi-relative-gap-after-obstruction-finite-resonant-subspace}
Let homology barrier induce closed subgroups \(\mathcal{H}\subset\widehat{\ZZ}\), and note that the twist transfer operator family is \(A(\chi)\)。
then for any role \(\chi\) satisfy \(\chi|_{\mathcal{H}}\neq 1\), there is a unified constant \(\delta>0\) make
\[
\boxed{\ \rho(A(\chi))\ \le\ (1-\delta)\rho(A(1))\ }.
\]
Therefore, even if there is a finite-mode obstacle, the uniform twisted spectral gap and relative improvement are still obtained on the orthogonal complement of the obstacle role-spanned subspace; the resonance is strictly restricted to a computable finite Abelian group, and its structure can be explicitly characterized by homologous torsion and Smith standard form.
\end{conclusion}
\begin{proof}[Points of proof]
Relative spectral gap and orthogonal complement theorem \ref{thm:gm-relative-gap-after-obstruction}；
See theorems for the homology characterization and Smith form computability of resonant role groups respectively. \ref{thm:gm-resonant-character-group-homology} with proposition \ref{prop:gm-smith-compute-obstructions}。
\end{proof}

\paragraph{\texorpdfstring{$p$}{p}-Congruential compression of power towers: Frobenius sparsification, Chebyshev closure and derivative filtering}

\begin{conclusion}[\texorpdfstring{$p^k$}{p^k}-Universal Frobenius law of sparsification: module \texorpdfstring{$p^\ell$}{p^ell} Next only "see"\texorpdfstring{$u^{p^{k-\ell+1}}$}{u^{p^{k-l+1}}}]\label{con:xi-witt-frobenius-sparsification-universal}
Periodic trace polynomial \(a_{p^k}(u)\in\ZZ[u]\), Dwork type congruence derived: in the quotient ring \((\ZZ/p^\ell\ZZ)[u]\) middle,\(a_{p^k}(u)\) depends only on variables \(u^{p^{k-\ell+1}}\), equivalently
\[
\boxed{\ a_{p^k}(u)\bmod p^\ell\in(\ZZ/p^\ell\ZZ)\bigl[u^{p^{k-\ell+1}}\bigr]\ }.
\]
therefore \(p\)-Power tower's high-level modular information can be deterministically pushed back to the low-level data, achieving structural compression where the modular audit complexity grows logarithmically with the level.
\end{conclusion}
\begin{proof}[Points of proof]
See proposition \ref{prop:witt-frobenius-iterated-descent}(Obtained by Dwork’s congruential step-by-step model reduction iteration).
\end{proof}

\begin{conclusion}[\texorpdfstring{$p=2$}{p=2} Chebyshev--Dwork closure: matrix power congruential recursive dimensionality reduction into single variable substitution \texorpdfstring{$t\mapsto t^2-2$}{t->t^2-2}]\label{con:xi-chebyshev-dwork-closure-p2}
at invariant coordinates \(t=u+u^{-1}\), even palindrome traces \(a_{2^k}(u)=u^{2^{k-1}}A_{2^k}(t)\) satisfies the congruential closed recursion
\[
\boxed{\ A_{2^k}(t)\equiv A_{2^{k-1}}(t^2-2)\pmod{2^k}\qquad(k\ge 2)\ }.
\]
Therefore, the fixed mold \(2^k\) can be calculated to convert an object from a polynomial matrix to a power \(B(u)^{2^k}\) is reduced to a logarithmic depth iteration of "single variable substitution + modulus"; this closure explicitly identifies the Frobenius push-forward on the two-power tower as a Chebyshev dynamical system.
\end{conclusion}
\begin{proof}[Points of proof]
See theorem \ref{thm:chebyshev-dwork-congruence-chain} and inference \ref{cor:chebyshev-dwork-matrixfree}。
\end{proof}

\begin{conclusion}[\texorpdfstring{$\theta$}{theta}-derivative \texorpdfstring{$p$}{p}-Advanced filtering: moments satisfy strict scaling congruence and lower bound of estimate]\label{con:xi-theta-derivative-padic-filter}
make \(u=e^\theta\), and define the derivative moment
\[
A_{n}^{(r)}:=\left.\frac{\mathrm{d}^r}{\mathrm{d}\theta^r}\,a_n(e^\theta)\right|_{\theta=0}\in\ZZ\qquad(r\ge 0).
\]
Then Dwork congruence implies step-by-step scaling congruence
\[
\boxed{\ A_{p^k}^{(r)}\equiv p^r\,A_{p^{k-1}}^{(r)}\pmod{p^k}\ }\qquad(r\ge 0),
\]
especially when \(0\le r<k\) when \(p\)-Enter the lower bound of valuation
\[
\boxed{\ p^{\,k-r}\ \big|\ A_{p^k}^{(r)}\ }.
\]
Therefore in \(p\)-power tower, all higher-order derivative moments automatically fall into strict \(p\)-into the filter window, giving a verifiable \(p\)-Enter rigid certificate.
\end{conclusion}
\begin{proof}[Points of proof]
See proposition \ref{prop:witt-theta-derivative-filter}。
\end{proof}

\paragraph{Natural boundaries and degree-of-freedom budget: non-extendibility constraints on the number of prime levels and zero-point density thresholds}

\begin{conclusion}[primitive--natural boundary of Dirichlet series:\texorpdfstring{$\Re(s)=0$}{Re(s)=0} Upper dense branch singularity enforces non-extendibility]\label{con:xi-primedirichlet-natural-boundary}
Dirichlet series for primitive orbitals \(\mathcal{P}(s)=\sum_{n\ge 1}p_n\lambda^{-ns}\), can be passed M\"obius--Abel type represents in \(\Re(s)>0\) analytic continuation;
But for every odd prime number \(p\) and any \(m\in\ZZ\),point
\[
s_{p,m}:=\frac{1}{p}+\frac{2\pi \mathrm{i} m}{p\log\lambda}
\]
is a logarithmic branch singularity, and \(\{s_{p,m}\}\) in a straight line \(\Re(s)=0\) is dense.
therefore \(\Re(s)=0\) is a natural boundary: there is no local analytic continuation across the straight line.
This constraint states that the "prime shadow" error term must be organized first by \(\zeta\) (or its logarithmic derivative/explicit formula) rather than a direct primitive--Dirichlet series to achieve controlled superposition.
\end{conclusion}
\begin{proof}[Points of proof]
See theorem \ref{thm:real-input-40-primedirichlet-dense-branch}(Dense branching singularity mechanism and natural boundary inference).
\end{proof}

\begin{conclusion}[Finite-dimensional Dirichletized \texorpdfstring{$\Theta(\log T)$}{Theta(log T)} Threshold: Fixed-bandwidth paradigm structurally fails at zero-point density order]\label{con:xi-finite-dimension-dirichlet-logT-threshold}
If the target is approaching \(\Xi_\zeta\) type zero counting main order \(\widetilde N_\zeta(T)\sim \tfrac{T}{\pi}\log T\), then any implementation that stays in "fixed finite bandwidth/fixed finite scale window/fixed finite dimension determinant aperture" for a long time will have a symmetric zero point count of at most \(O(T)\), thus inevitably failing at the density level.
Therefore it must be released \(\Theta(\log T)\)-level degree-of-freedom budget (e.g. using an exponential \(\tau(T)\) follow \(T\) grow and satisfy \(\tau(T)\gtrsim \log T\), into a growing dimensional approximation sequence or an infinite dimensional Fredholm determinant caliber).
\end{conclusion}
\begin{proof}[Points of proof]
The Cartwright density upper bound is given by \(\widetilde N_T(T)\le \frac{2}{\pi}\tau(T)T+o(\tau(T)T)\)；
To match \(\widetilde N_\zeta(T)\sim \frac{T}{\pi}\log T\), must have \(\tau(T)\gtrsim \frac12\log T\)。
See conclusion \ref{con:xi-terminal-zero-density-product-lower-bound}。
\end{proof}

\paragraph{Phase Closure and Unit Circle Root Certificate Chains: Fixed Slices, Jensen Defects, and Schur--Cohn/LMI}

\begin{conclusion}[Unique fixed slice of phasing closure: zero point is institutionally locked at when external radial registers are rejected \texorpdfstring{$\Re(s)=\tfrac12$}{Re(s)=1/2}]\label{con:xi-unitary-slice-unique-fixed-slice}
In the unitary-slice readout system, if the verification object can be expressed as
\(
D(s)=\det(I-u_L(s)U)
\)
（\(U\) is the unitary operator,\(u_L(s)=L^{2s-1}\)), and the system rejects external radial/amplitude registers on the mode axis, then any addressable zero point event must satisfy \(|u_L(s)|=1\),thereby \(\Re(s)=\tfrac12\)。
Equivalently,\(\Re(s)=\tfrac12\) is the only fixed slice of the phased closure in the completion coordinates, which is of finite type Hilbert--P{\'o}The lya route provides a strong binding framework within the system.
\end{conclusion}
\begin{proof}[Points of proof]
Depend on \(u_L(s)=L^{2s-1}\) have to \(\Re(s)=\tfrac12\Longleftrightarrow |u_L(s)|=1\)；
Take from visible domain \(\mathsf{Phase}\) and reject the external radial register,\(|u_L(s)|\neq 1\)'s click protocol is read as \(\mathsf{NULL}\), so that the addressable zero point is locked to the critical line.
See \S\ref{subsec:xi-pw-type-safety-null}(especially the theorem \ref{thm:xi-offset-null-type-safety}）。
\end{proof}

\begin{conclusion}[Jensen defect as horizon entropy distance: deviations can be continuously quantified and verified]\label{con:xi-jensen-defect-horizon-entropy-distance}
Complete the self-inversive polynomial \(P\)(satisfy \(y^dP(1/y)=e^{\mathrm{i}\phi}P(y)\), and assume \(P(0)\neq 0\)),right \(0<\varrho<1\) defines the Jensen defect
\[
D_P(\varrho):=\frac{1}{2\pi}\int_{0}^{2\pi}\log|P(\varrho e^{\mathrm{i}\theta})|\,\dd\theta-\log|P(0)|.
\]
Then "all zero points are on the unit circle" and \(D_P(\varrho)\equiv 0\)（\(0<\varrho<1\)) is equivalent to
\(\lim_{\varrho\uparrow 1}D_P(\varrho)=0\)。
therefore \(D_P(\varrho)\) can be interpreted as a continuous defect function for radial/amplitude leakage: while holding the self-inversive (completed self-dual) structure constant, it gives an auditable distance scale from "strict Ramanujan/RH-sharp" to "substantial defect".
\end{conclusion}
\begin{proof}[Points of proof]
right \(P\) applying Jensen’s formula, we get \(D_P(\varrho)=\sum_{|a_k|<\varrho}\log(\varrho/|a_k|)\ge 0\)；
therefore \(D_P(\varrho)\equiv 0\) is equivalent to any \(\varrho<1\) There is no internal zero point.
like \(P\) self-inversive, then the zero point is about \(y\mapsto 1/\overline y\) are in pairs, so no zero point in the disk is equivalent to all zero points falling on the unit circle.
For the functional caliber of this manuscript, see Proposition \ref{prop:app-jensen-formula-defect} and theorem \ref{thm:app-rh-iff-disk-zero-free}。
\end{proof}

\begin{conclusion}[Cayley hyperhyperbolicization and Cohn derivative intradiscretion: the unit circle root criterion compressed into a compilable positive semidefinite certificate chain]\label{con:xi-cayley-cohn-lmi-certificate-chain}
For self-reciprocal real coefficient polynomials \(\Phi_L\) can be equivalently transformed into a real polynomial of one variable through Cayley transformation \(\Psi_L\) of all real roots (hyperhyperbolicity);
At the same time, according to Cohn's criterion, the complete root of the unit circle is equivalent to \(\Phi_L\) self-inversive and \(\Phi_L'\) all roots fall within the closed unit disk.
Therefore, the unit circle root geometry can be compressed into two achievable audit paths: the real-axis hyperbolicity path and the derivative indiskness path; both can be further connected to the Schur--Cohn/Toeplitz type positive semi-definite certificate, thereby institutionalizing the root position assertion into a verifiable LMI chain.
\end{conclusion}
\begin{proof}[Points of proof]
Cayley's hyperbolic solution proposition \ref{prop:xi-rh-cayley-hyperbolicity}；
See the proposition for the intra-discrimination of Cohn type derivatives \ref{prop:xi-rh-cohn-derivative}；
See Notes for the alignment of the Schur--Cohn positive semidefinite interface with the Toeplitz certificate. \ref{rem:xi-schur-cohn-sdp-interface} and theorem \ref{thm:xi-toeplitz-schur-cohn-lift}。
\end{proof}

\paragraph{Detectable spectral fingerprints and transient depth of limited samples: Edgeworth constant, nilpotent exponent and Witt--Dwork inheritance}

\begin{conclusion}[Third-order spectral fingerprint of gauge difference fluctuations: Edgeworth coefficient gives universal negative skewness constant and locks Jordan defect]\label{con:xi-gauge-anomaly-third-order-spectral-fingerprint}
standard difference process \(G_m\) central limit theorem scale normalization of
\[
Z_m:=\frac{G_m-\frac49 m}{\sqrt{\frac{118}{243}m}},
\]
Its first-order Edgeworth correction term coefficient has a completely closed form
\[
\boxed{\;
\gamma_1^{(G)}
=\frac{P_G^{(3)}(0)}{(P_G''(0))^{3/2}}
=\frac{-\frac{1174}{2187}}{\left(\frac{118}{243}\right)^{3/2}}\;<\;0\;}
\]
The skewness is thus strictly negative and uniquely locked by finite-dimensional spectral/Jordan data.
The third-order constant turns \((-1/2)\)'s Jordan defect is upgraded from a structural narrative to a numerical fingerprint that can be detected with a limited sample.
\end{conclusion}
\begin{proof}[Points of proof]
The closed form of the third-order pressure derivative is shown in the proposition \ref{prop:fold-gauge-anomaly-pressure}；
Edgeworth first-order sealing and coefficient expression see corollary \ref{cor:fold-gauge-anomaly-edgeworth1}。
\end{proof}

\begin{conclusion}[The nilpotent index is the transient memory depth: the kernel dimension chain provides a strict upper bound and is the same as the length of the reset word]\label{con:xi-nilpotent-index-transient-memory-depth}
to real input \(40\) full matrix of state kernel \(M\), kernel dimension chain \(\dim\ker(M^k)\) is given by the stability point of \(0\) Maximum Jordan block size at eigenvalues \(r_{\mathrm{full}}=5\)；
Therefore for any fixed vector \(v,w\),sequence \(f_m=v^\ast M^m w\) exist \(m\ge r_{\mathrm{full}}\) followed by \(w\) exist \(0\) - component independent on the transient subspace of the spectrum (the transient contribution is completely annihilated by the nilpotent part).
In the same way, the corresponding index of the essential kernel is \(3\)。
This invariant improves the online delay and reset word length to hard upper bounds that can be recalculated by a linear algebra chain with one click.
\end{conclusion}
\begin{proof}[Points of proof]
See proposition \ref{prop:real-input-40-nilpotent-index} and inference \ref{cor:real-input-40-nilpotent-memory-depth}(For reproducible output of kernel dimension chains and Jordan block counts see note \ref{rem:real-input-40-nilpotent-index-repro}）。
\end{proof}

\begin{conclusion}[Zero-cost inheritance of Big-Witt/Dwork guardrails: multi-variable upgrades do not change the divisibility and congruence structures]\label{con:xi-big-witt-dwork-guardrail-multivariable}
2D/3D weighting matrix \(B(u,v)\)、\(B(u,v,w)\) trace polynomial \(a_n\) still satisfies the divisibility law of necklace and the literal generalization of Dwork-type congruence:
\[
\sum_{d\mid n}\mu(d)\,a_{n/d}(u^d,v^d)\in n\ZZ[u,v],
\qquad
a_{p^k}(u,v)\equiv a_{p^{k-1}}(u^p,v^p)\pmod{p^k},
\]
The same goes for the three-dimensional case.
Therefore the primitive--periodic Witt Euler product is \(p\)-Power congruence remains unchanged under multidimensional energy resolution, providing an implementation self-auditing guardrail: the divisibility/congruence consistency check can be completed simply by tracing the trace sequence, without re-entering the costly domain of matrix power calculations.
\end{conclusion}
\begin{proof}[Points of proof]
See the conclusion for the univariate template of Witt Euler product and Dwork congruence \ref{con:xi-terminal-big-witt-functoriality} and inference \ref{cor:witt-dwork-congruence}；
Multivariable generalization is directly deduced from the functor property of "variable substitution to maintain the congruence of integral coefficients and modules" and the item-by-item definition of ghost coordinates (also consistent with \S\ref{subsec:operator-cyclic-traceclass-witt-rigidity} The trace sequence has the same caliber).
\end{proof}



% (User input window)
%
% (Source tag) /killo-golden
% Timestamp: 2026-02-18 14:26:17 (Asia/Singapore)
%
% The content is rewritten using the existing notation and auditable interface of this article, without introducing manual numbering.

\paragraph{Stable negative subspaces, endpoint spectra and temporal fibers: a compilable interface for Toeplitz falsifications}
This section further compresses the negative spectrum structure of the "finite falsifier" into a recalculable model space truncated image and endpoint readout invariant, and transforms the window-\(6\)'s minimal thermal closure, Witt--Fourier double bridge and synchronized nuclear center sixfold aggregation are incorporated into the same set of automatically numbered conclusion calibers.

\begin{conclusion}[Compilability of Toeplitz denier: Stable negative subspace is a truncated image of the model space]\label{con:xi-toeplitz-compiled-model-space-truncation}
Under the limited Blaschke defect setting, let \(\kappa:=\deg(B_{F_{\zeta}})=\dim K_{B_{F_{\zeta}}}\)。
then exists \(N_0\) follow the clan \(N\) Stable explicit linear embedding
\[
\iota_N:K_{B_{F_{\zeta}}}\hookrightarrow \CC^{N+1}\qquad(N\ge N_0)
\]
(For example, taking the Hardy basis \(\{1,w,w^2,\dots\}\) before \(N{+}1\) Taylor coefficient coordinates), so that the Toeplitz quadratic form
\(
x\mapsto x^\ast\mathbf{T}_N x
\)
exist \(\iota_N(K_{B_{F_{\zeta}}})\) has a limit on\textbf{exactly \(\kappa\) Dimensional negative definite part}, and the negative definite part is in \(N\to\infty\) remains unchanged within the stable interval.
So once it is observed \(\mathbf{T}_N\), we can realize a compression shift with \(A_{B_{F_\zeta}}\)'s point spectrum residual axis is read out with the same type of "residual axis", and its point spectrum data is used to invert the defect points (in terms of multiplicity) in the disk.
\end{conclusion}
\begin{proof}[Points of proof]
Stable embedding and "negative subspace \(=\) See the theorem for the isomorphism of “truncated image of model space” \ref{thm:xi-toeplitz-witness-generated-by-model-space}；
The correspondence between the point spectrum residual axis and the Blaschke zero point is shown in the theorem \ref{thm:blaschke-point-spectrum-correspondence}。
Combined with spectral readout interfaces (e.g. theorem \ref{thm:xi-resonance-pencil-reconstruction} and conclusion \ref{con:xi-finite-falsifier-offline-fiber-readout}) can upgrade the stable negative subspace to a finite-dimensional auditable defect location inversion.
\end{proof}

\begin{conclusion}[Minimal generative law of failed witnesses: falsification is forced by the point spectrum residual axis]\label{con:xi-toeplitz-witness-minimal-generator-axis}
The limited defect setting based on the previous conclusion is used.
when \(N\) is sufficiently large, all negative directions of the Toeplitz level can be determined by the point spectrum residual axis
\(
\mathcal H_{\mathrm{pp}}(A_{B_{F_\zeta}})=K_{B_{F_\zeta}}
\)
point mode of \(\iota_N\); thus the "minimum generating dimension" of Toeplitz's falsification is strictly equal to
\(
\deg(B_{F_\zeta})=\kappa
\)。
Finite falsifiers are therefore auditable finite-dimensional generated objects rather than diffuse by-products of numerical instability.
\end{conclusion}
\begin{proof}[Points of proof]
See the theorem for the stability of negative inertia index on high-order truncation \ref{thm:xi-toeplitz-inertia-stabilizes-to-kappa}；
Point mode generativity and minimum dimension result from the theorem \ref{thm:xi-toeplitz-witness-generated-by-model-space} given directly.
\end{proof}

\begin{conclusion}[Identical operatorization of endpoint flux: Poisson--Busemann weight is equal to the trace of Julia index and endpoint evaluation operator]\label{con:xi-endpoint-flux-operator-trace-identity}
set up \(B\) is the finite Blaschke product, and the zero-point multiset is \(\{a_k\}_{k=1}^{\kappa}\subset\mathbb{D}\), and note the endpoint Poisson weight
\[
P_a(-1)=\frac{1-|a|^2}{|1+a|^2}.
\]
Define endpoint flux (absorption coefficient)
\[
\Phi_{-1}(B):=\sum_{k=1}^{\kappa}P_{a_k}(-1),
\]
And order \(K_B:=H^2\ominus BH^2\), endpoint evaluation functional \(L_{-1}(f):=f(-1)\)。
Then there is a strict identity
\[
\boxed{\ \Phi_{-1}(B)=\mathcal{J}_{-1}(B)=\mathrm{tr}\!\bigl(L_{-1}^\ast L_{-1}\bigr)\ },
\]
As a result, the endpoint absorption intensity is promoted to the spectrozable trace invariant of the "endpoint evaluation operator", and is completely equivalent to the point mode boundary coupling sum of squares.
\end{conclusion}
\begin{proof}[Points of proof]
This is the theorem \ref{thm:xi-endpoint-flux-pointmode-coupling} conclusion (and use \eqref{eq:app-poisson-kernel-endpoint-minus1} The Poisson endpoint is closed and interfaces with Julia indicators).
\end{proof}

\begin{conclusion}[Endogeneity of time fibers: the imaginary part of the upper half plane is the endpoint \texorpdfstring{$-1$}{-1} Poisson--Busemann density]\label{con:xi-timefiber-poisson-busemann-identity}
make \(t=x+\mathrm{i}y\in\mathbb{H}\), and take the Cayley transform
\[
w=\mathcal{C}(t)=\frac{1+\mathrm{i}t}{1-\mathrm{i}t}\in\mathbb{D}.
\]
then a strict identity is established
\[
\boxed{\ y=\mathcal{B}_{-1}(w)=\frac{1-|w|^2}{|1+w|^2}\ }.
\]
Therefore, under the endpoint compaction caliber, the "time fiber coordinates" are not additional continuous axes, but are strictly identical to the endpoints. \(-1\) of the harmonic measure density (Poisson--Busemann weight); high-height information is forced to be compressed into \(w\approx-1\), thus forming a verifiable endpoint localization principle.
\end{conclusion}
\begin{proof}[Points of proof]
See corollary \ref{cor:dephys-poisson-busemann-timefiber}(and cite the proposition \ref{prop:app-busemann-poisson-minus1} Busemann/Poisson defined interface).
\end{proof}

\paragraph{window-\texorpdfstring{$6$}{6} Minimal thermal closure of : Push kernels, spectral gaps and lumpability barriers}

\begin{conclusion}[window-\texorpdfstring{$6$}{6} Minimal thermal closure of : Reversibility, detailed equilibrium and explicit spectral gaps of push kernels]\label{con:xi-window6-p6-minimal-thermalization-closure}
Push adjacency count \(A_6\) definition \(X_6\) on the push core \(P_6\) and stationary distribution \(\pi_6\) (Generate fragment \eqref{eq:window6_pushforward_markov_kernel}）：
\[
\pi_6(x)=\frac{d^{\mathrm{bin}}_6(x)}{64},\qquad
P_6(x,x')=\frac{A_6(x,x')}{6\,d^{\mathrm{bin}}_6(x)}.
\]
but \(P_6\) satisfies delicate balance
\(
\pi_6(x)P_6(x,x')=\pi_6(x')P_6(x',x)
\)，
Therefore, it is reversible and \(\pi_6\) is a stationary distribution.
Spectral gap auditing gives auditable constants (generated fragments \eqref{eq:window6_pushforward_markov_spectral_gap}）
\[
\lambda_2(P_6)\approx 0.4841207858,\qquad
\lambda_\star(P_6)\approx 0.6030939754,
\]
So for any initial distribution \(\mu\) exists constant \(C(\mu)\) applies to all \(t\ge 0\) has an exponential mixture upper bound
\[
\|\mu P_6^t-\pi_6\|_{\mathrm{TV}}\le C(\mu)\,\lambda_\star(P_6)^t.
\]
\end{conclusion}
\begin{proof}[Points of proof]
Kernel and detailed balance see theorem \ref{thm:terminal-foldbin6-pushforward-markov} and \eqref{eq:window6_pushforward_markov_kernel}；
Spectral gap constant and TV upper bound see corollary \ref{cor:terminal-foldbin6-minimal-markov-closure-gap} and \eqref{eq:window6_pushforward_markov_spectral_gap}。
\end{proof}

\begin{conclusion}[Hard constant from spectral gap to mixing time: window-\texorpdfstring{$6$}{6} The thermalization time scale of \texorpdfstring{$1/(-\log\lambda_\star)$}{1/(-log lambda*)}]\label{con:xi-window6-mixing-time-hard-constant}
Use the notation for the previous conclusion.
If for a given initial distribution \(\mu\) defines the convergence time
\[
t(\varepsilon;\mu):=\min\bigl\{t\ge 0:\ \|\mu P_6^t-\pi_6\|_{\mathrm{TV}}\le \varepsilon\bigr\},
\]
then by \(\|\mu P_6^t-\pi_6\|_{\mathrm{TV}}\le C(\mu)\lambda_\star(P_6)^t\) to get an explicit upper bound
\[
\boxed{\ t(\varepsilon;\mu)\ \le\ \frac{\log(C(\mu)/\varepsilon)}{-\log\lambda_\star(P_6)}\ }
\approx 1.9775264\cdot \log\!\bigl(C(\mu)/\varepsilon\bigr).
\]
Therefore window-\(6\) carries an implementation-independent protocol-level hard constant on the time scale \(1/(-\log\lambda_\star(P_6))\), and are output directly from the spectral certificate without relying on additional fitting.
\end{conclusion}
\begin{proof}[Points of proof]
by inference \ref{cor:terminal-foldbin6-minimal-markov-closure-gap} The TV exponential convergence formula of is obtained by taking the logarithm;
The numerical constant is given by \eqref{eq:window6_pushforward_markov_spectral_gap} in \(\lambda_\star(P_6)\) is substituted into the calculation.
\end{proof}

\begin{conclusion}[Structural consequences of strong merging failure: Stable type processes cannot be first-order Markov closed in general initial states]\label{con:xi-window6-strong-lumpability-failure-structural}
divided by fiber \(\{F_6^{-1}(x)\}_{x\in X_6}\)-induced coarse-graining does not satisfy strong lumpability:
There are two points in the same fiber \(\omega,\omega'\) makes its one-step neighbors have different count distributions on stable types (for an explicit counterexample see \eqref{eq:foldbin6_strong_lumpability_counterexample}）。
Therefore there is no dependency on anything other than \(X_t\) homogeneous transfer kernel \(P\) can be in\emph{all}Under the initial microstate distribution, make \((X_t)\) is closed as \(X_6\) on a memoryless Markov chain.
\end{conclusion}
\begin{proof}[Points of proof]
This is a proposition \ref{prop:terminal-foldbin6-strong-lumpability-fails};Counterexample generated by fragment \eqref{eq:foldbin6_strong_lumpability_counterexample} given.
\end{proof}

\begin{conclusion}[Forcing of minimal memoryless closure: instantaneous thermalization within fibers is strictly equivalent to push core closure]\label{con:xi-window6-instant-thermalization-equivalence}
In the context of strong compilability failure, if a minimal memoryless closure mechanism is introduced: after each step of cube flipping, the condition is the current stable type \(X_t=x\) immediately homogenizes the microstate to the fiber \(F_6^{-1}(x)\) (instantaneous heating within the fiber),
Then the stable type process is closed as \(X_6\), its transfer core must be a push core \(P_6\), the equilibrium measure must be \(\pi_6\propto d^{\mathrm{bin}}_6\)。
Equivalently, under this protocol caliber, the existence of first-order memory-less closure is strictly equivalent to "instantaneous thermalization within the fiber".
\end{conclusion}
\begin{proof}[Points of proof]
The uniqueness and reversibility of closure are given by the theorem \ref{thm:terminal-foldbin6-pushforward-markov} and inference \ref{cor:terminal-foldbin6-minimal-markov-closure-gap} is given; its compulsion is given by the proposition \ref{prop:terminal-foldbin6-strong-lumpability-fails} The negative conclusion is derived.
\end{proof}

\paragraph{Witt--Fourier Double Bridge and Operatorization Big-Witt: Unifying the Ghost Ruler and \texorpdfstring{$p$}{p}-Typical guardrail}

\begin{conclusion}[Witt--Fourier The strict bonding of two bridges: the same character axis \texorpdfstring{$\mathrm{Tw}_\chi$}{Tw\_chi} Simultaneously diagonalize "convolution--collision" and "trace--Euler"]\label{con:xi-witt-fourier-commuting-cube}
in finite ring \(G_m\) and its character group \(\widehat G_m\), for each role \(\chi\) defines fiber side twist
\[
\mathrm{Tw}_\chi(d_m):=\sum_{r\in G_m}d_m(r)\chi(r)=\widehat d_m(\chi),
\]
And on the track side with the same \(\chi\) defines the Dirichlet twist kernel \(T_\chi\)。
Then the Fourier diagonalization and the Witt/Euler exchange diagram are on the same line \(\mathrm{Tw}_\chi\) are strictly compatible on the axis, and the two are bonded along the character direction to form a strict exchange cube;
And there is complete factorization under the Fold caliber
\[
\mathrm{Tw}_\chi(d_m)=\prod_{j=1}^m\bigl(1+\chi(F_{j+1})\bigr).
\]
\end{conclusion}
\begin{proof}[Points of proof]
See note for consistency between swapping cubes and twisting axes \ref{rem:twist-commuting-cube}；
See the theorem for multiplication and integral solution \ref{thm:fold-spectrum-factorization}(And give the alignment in the same note).
\end{proof}

\begin{conclusion}[Unified ghost ruler: the second-order collision gap simultaneously measures the deviation intensity of Fourier non-trivial spectral energy and Witt/Euler]\label{con:xi-collision-gap-ghost-yardstick}
The "trace" on Witt's side \(\to\) primitive \(\to\) Euler" bridge package and the Parseval energy bridge package on the Fourier side share the same computable scalar fingerprint - the second-order collision gap (equivalent to the non-trivial spectral energy \(\ell^2\)-norm):
\[
\boxed{\ F_{m+2}\,\mathsf{Col}_m-1=\frac{1}{4^m}\sum_{t\ne 0}\bigl|\widehat c_m(t)\bigr|^2\ }.
\]
Therefore, the deviation intensity of the prime-like (Euler product) side and the spectrum/trace side can be compressed to the same auditable scalar, forming the shortest verification loop for cross-channel consistency verification.
\end{conclusion}
\begin{proof}[Points of proof]
See note \ref{rem:fw-bridge}(And by the theorem \ref{thm:fold-collision-spectrum} The Parseval identity of gives a rigorous derivation of this scale; its equivalence to Rényi-2/uniformity gap is shown in the corollary \ref{cor:pom-renyi2-near-uniform}）。
\end{proof}

\begin{conclusion}[Operator compilation of Big-Witt operations: semi-ring homomorphically embedded trace operator combination of atomic Euler data]\label{con:xi-atomic-witt-operator-compilation}
make \(\mathcal{W}_{\mathrm{atom}}\) is the atomic data \((\beta,n,m)\) generates a Big-Witt type semiring, let \(\mathrm{TC}_1\) is the commutative semiring of the trace operator under the unitary equivalence class (addition is a direct sum and multiplication is a tensor product).
Then the mapping
\[
(\beta,n,m)\longmapsto \bigoplus_{i=1}^{m} C(n,\alpha),\qquad \alpha^n=\beta,
\]
Extension to constructive semicyclic homomorphism
\[
\boxed{\ \mathcal{W}_{\mathrm{atom}}\ \longrightarrow\ \mathrm{TC}_1\ }.
\]
Therefore Big-Witt's abstract Euler/plethystic operations can be compiled into direct sum/tensor product combination operations of trace operators, and combined with Fredholm \(\zeta\) interface is strictly consistent, forming an executable Euler algebra.
\end{conclusion}
\begin{proof}[Points of proof]
See theorem \ref{thm:atomic-witt-into-tc1}(multiplicative \(\gcd/\mathrm{lcm}\) rules consist of propositions \ref{prop:cyclic-block-tensor-gcd-lcm} ensure).
\end{proof}

\begin{conclusion}[\texorpdfstring{$p$}{p}-Typical Frobenius--Verschiebung's trace layer implementation:\texorpdfstring{$\mathsf F_p\circ\mathsf V_p$}{Fp o Vp} Equivalent to "direct sum times \texorpdfstring{$p$}{p}”]\label{con:xi-ptypical-frobenius-verschiebung-trace}
exist \(p\)-on typical components, the operator implementations of Frobenius and Verschiebung satisfy ghost compatibility and the Verschiebung--Frobenius relationship:
for any reason \(p\)-Operator obtained by finite direct sum of power loop blocks \(L\), to everything \(n\ge 1\) have
\[
\boxed{\ \Tr\!\bigl(\mathsf F_p(L)^n\bigr)=\Tr\!\bigl(L^{pn}\bigr)\ },
\qquad
\boxed{\ \mathsf F_p\!\bigl(\mathsf V_p(L)\bigr)\ \simeq\ L^{\oplus p}\ }.
\]
This relationship promotes the Dwork type congruence from "polynomial/sequence" to the structural identity of "trace operator combination", making \(p\)-Enter guardrails are inherited at zero cost in the way of operator functors.
\end{conclusion}
\begin{proof}[Points of proof]
See proposition \ref{prop:cyclic-p-typical-frobenius-verschiebung}。
\end{proof}

\paragraph{Modular barriers, cyclotomic factors, and Smith shapes: a finite termination audit of resonance role tables}

\begin{conclusion}[Complete equivalence of modular barriers and cyclotomic factors: Resonance sources reduced to unit root factor classification]\label{con:xi-mod-obstruction-cyclotomic-equivalence}
to ihara/kinetics \(\zeta\) of the twisted channel \(\det(I-tB_\rho(u))^{-1}\), there exists a non-trivial modular obstacle if and only if some \(\rho\neq \mathbf{1}\) Make the corresponding spectral radius reach the trivial channel;
And this condition is if and only if \(\det(I-tB_\rho(u))\) contains cyclotomic factors (unit root factors appear on the unit circle after normalization).
Therefore, the "resonance source" can be equivalently reduced to the cyclotomic factor classification problem of finite-dimensional integral coefficient polynomials.
\end{conclusion}
\begin{proof}[Points of proof]
See conclusion \ref{con:xi-mod-obstruction-cyclotomic-factor} and its cited classification theorem \ref{thm:conclusion75-ihara-cyclotomic-factor-classification}。
\end{proof}

\begin{conclusion}[Cohomology classification and polynomial-time computability of obstacle groups: a fully resonant role table can be compiled into a single SNF calculation]\label{con:xi-resonant-character-snf-computable}
Resonance character collection \(X_{\mathrm{res}}\) constitutes a finite Abelian group and can be characterized by the torsion part of the homology quotient;
Its invariant factorization and generation roles can be read directly from the Smith normal form of the corresponding integer matrix.
More strongly, there are deterministic algorithms that work in time \(\mathrm{poly}(|V|,|E|,\log W)\) inner output barrier invariant and \(X_{\mathrm{res}}\) and generate roles, so that the full-role audit of "whether there are cyclotomic factors/resonance channels" can be terminated in finite steps.
\end{conclusion}
\begin{proof}[Points of proof]
See the theorem for homology characterization \ref{thm:gm-resonant-character-group-homology}; Smith form and polynomial time computability see the proposition \ref{prop:gm-smith-compute-obstructions}。
\end{proof}

\paragraph{Central sixfold aggregation and six-branch inertia of the Perron branch: analytical radius and algebraic degree constraints}

\begin{conclusion}[Geometric consequences of central sixfold aggregation: The distance between the unique real intersection of the center and the nearest complex singular point gives the analytical radius]\label{con:xi-rate-center-sixfold-unique-real-rstar}
Normalized Elimination Polynomials for Synchronous Kernel Rate Curves \(R(\alpha,u)\) has a dual self-symmetric structure (Note \ref{rem:rate-algebraic-elimination-structure}), and satisfies the decomposition at the central section
\[
R\!\left(\tfrac12,u\right)=-\frac{(u-1)^6}{64}\,Q(u).
\]
in \(Q(u)\) coefficient palindrome, so \(Q(u)=u^{16}Q(1/u)\), zeros appear in inverse pairs (and merge with complex conjugates into quadruples); and Sturm counting gives \(Q\) has no zero point on the real axis, so \(Q(u)>0\) for all \(u\in\RR\) is established (proposition \ref{prop:rate-center-Q-reciprocal-no-real}）。
therefore \(R(\tfrac12,u)=0\) on real parameters if and only if \(u=1\), and the multiplicity is exactly \(6\)(inference \ref{cor:rate-center-unique-real}）。
Further define the nearest complex singular point distance
\[
R_\star:=\min\{|u-1|:\ Q(u)=0\},
\]
but \(R_\star\) gives the analytic radius of the central branch and can be numerically audited (corollary \ref{cor:rate-center-Rstar}）。
\end{conclusion}
\begin{proof}[Points of proof]
See note for recomputable certificates of self-symmetry and central decomposition \ref{rem:rate-algebraic-elimination-structure}；
\(Q\)'s self-reciprocity and real root emptiness see the proposition \ref{prop:rate-center-Q-reciprocal-no-real}；
The corollary of the unique true intersection of the center and the sixfold nature \ref{cor:rate-center-unique-real}；
For a numerical certificate of the distance to the nearest complex singularity, see Corollary \ref{cor:rate-center-Rstar}。
\end{proof}

\begin{conclusion}["Six-branch unit" forced by inertial group: Perron branch algebraic degree must be \texorpdfstring{$6$}{6} divisible]\label{con:xi-rate-center-perron-degree-multiple-6}
Follow the central six-fold aggregation certificate \(R(\tfrac12,u)=-(u-1)^6Q(u)/64\), and assuming general transversality conditions \(\partial_\alpha R(\tfrac12,1)\neq 0\)。
make \(F(\alpha,u)\in\QQ[\alpha,u]\) for \(R(\alpha,u)\) exist \(\QQ(\alpha)[u]\) contains points \((\tfrac12,1)\) (such as the one containing the Perron branch).
then expands in the function domain \(\QQ(\alpha)\subset \QQ(\alpha,u)\ (F(\alpha,u)=0)\) middle,\(\alpha=\tfrac12\), the divergence index is \(e=6\),thereby \(\deg_uF\) (that is, the branch pair \(\QQ(\alpha)\) must be \(6\) divisible.
Equivalently, the Perron branch of the central point neighborhood has an irreducible six-branch inertial structure, and any scheme that attempts to "cross the central point" with a single-valued analytic branch is ruled out at the algebraic level.
\end{conclusion}
\begin{proof}[Points of proof]
See corollary \ref{cor:rate-center-perron-degree-multiple-6};See notes for semantic alignment of inertia groups \ref{rem:rate-center-inertia-c6}。
\end{proof}

\endinput


% (User input window)
%
% (Source tag) /killo-golden
% Timestamp: 2026-02-18 14:57:09 (Asia/Singapore)
% Source: main.pdf (accessed on 2026-02-18)
%
% The content is rewritten using the existing notation and auditable interface of this article, without introducing manual numbering.

\paragraph{Negative spectral volume potential, stretching translation and two-state degeneration: Pick--Poisson invariant, Hankel homology pathological condition and recomputable DP kernel}
This section organizes the complete invariants of the Pick--Poisson negative spectral volume potential, the \(\log m\) translation under unified scaling, the homology pathological conditions of the two types of Hankel/Toeplitz inversions, as well as the discrete selection and two-state closure on the window-\(6\)/bin-fold side, organized into a set of items that can be directly cited; the entrance to their proofs can be accurately located by the existing theorems/propositions of this article.

\begin{conclusion}[Complete invariants of negative spectral determinants: Poisson weights and pseudohyperbolic Vandermonde product limits]\label{con:xi-terminal-negative-spectrum-det-invariant}
Let point defects be given by the finite Blaschke zero-point multiset \(\{a_j\}_{j=1}^{\kappa}\subset\mathbb{D}\), and let \(\mathsf P\) be the corresponding Pick--Poisson Gram matrix \eqref{eq:xi-pick-poisson-gram}.
Let pseudo-hyperbolic distance
\[
\rho(a,b):=\left|\frac{a-b}{1-\overline b\,a}\right|\qquad(a,b\in\mathbb{D}).
\]
Then the determinant has a closed form decomposition
\[
\boxed{\;
\det\mathsf P
=
\Bigl(\prod_{j=1}^{\kappa}P_{a_j}(-1)\Bigr)
\Bigl(\prod_{1\le j<k\le\kappa}\rho(a_j,a_k)^2\Bigr)\; }.
\]
And when the Toeplitz master submatrix \(\mathbf{T}_N\) enters the stable development interval and its negative eigenvalue (by multiplicity) is \(\{\lambda_j^{-}(\mathbf{T}_N)\}_{j=1}^{\kappa}\), there is a strict multiplication limit
\[
\boxed{\;
\lim_{N\to\infty}\ \prod_{j=1}^{\kappa}\bigl(-\lambda_j^{-}(\mathbf{T}_N)\bigr)=\det\mathsf P\; }.
\]
This invariant undergoes strict degeneration for both edge (\(P_{a_j}(-1)\downarrow 0\)) and collision (\(\rho(a_j,a_k)\downarrow 0\)), thus completely encoding the point-like defect geometry in the semi-classical limit sense.
\end{conclusion}
\begin{proof}[Proof sketch]
See proposition \ref{prop:xi-pick-poisson-det-factorization}.
\end{proof}

\begin{conclusion}[Covariant law of unified scaling: development complexity occurs \texorpdfstring{$\log m$}{log m} translation]\label{con:xi-terminal-visibility-complexity-logm-translation}
Use the finite defect setting of the previous conclusion, and let \(\phi_m(w):=\mathcal{C}(m\,\mathcal{C}^{-1}(w))\) (\(m>0\), fixed endpoint \(-1\)) to be uniform scaling.
Let \(a_j(m):=\phi_m(a_j)\), and let \(\mathsf P(m)\) be the Pick--Poisson Gram matrix generated by \(\{a_j(m)\}\) via \eqref{eq:xi-pick-poisson-gram}.
Define boundary entropy potential, boundary volume potential and development complexity
\[
S_{\partial}:=-\log\det\mathsf P,
\qquad
V_{\partial}:=\log\mathrm{tr}(\mathsf P),
\qquad
\mathscr{C}_{\partial}:=S_{\partial}+(\kappa-1)V_{\partial},
\]
And define \(S_{\partial}(m),V_{\partial}(m),\mathscr{C}_{\partial}(m)\) with the same formula.
Then strict covariance is established for all \(m>0\)
\[
\boxed{\;
\mathscr{C}_{\partial}(m)=\mathscr{C}_{\partial}(1)+\log m\; }.
\]
The scale parameter is therefore linearized into additive coordinates on \(\mathscr{C}_{\partial}\); this translation is mechanistically aligned with the translation identity \(y(ms)=y(s)+\log m\) (theorem \ref{thm:koenigs-linearization-witt-dilation}) of Witt dilation on the Koenigs depth coordinate \(y\).
\end{conclusion}
\begin{proof}[Proof sketch]
By the proposition \ref{prop:xi-pick-poisson-det-scaling-law}, \(\rho\) is invariant under \(\phi_m\) and \(P_{\phi_m(a)}(-1)=m^{-1}P_a(-1)\), so \(\det\mathsf P(m)=m^{-\kappa}\det\mathsf P\) and \(\mathrm{tr}(\mathsf P(m))=m^{-1}\mathrm{tr}(\mathsf P)\).
Substitute \(\mathscr{C}_{\partial}=S_{\partial}+(\kappa-1)V_{\partial}\) to get the conclusion; see the inference \ref{cor:xi-toeplitz-visibility-complexity-anomaly} for details.
\end{proof}

\begin{conclusion}[Electrostatic characterization of negative spectral volume potential:\texorpdfstring{$-\log\det\mathsf P$}{-log det P} is equivalent to the upper half-plane Dirichlet Green mirror energy]\label{con:xi-terminal-pick-poisson-green-mirror-energy}
Under the upper half-plane coordinate \(z_j=\mathcal{C}^{-1}(a_j)\in\mathbb{H}\), let
\(
t_j:=-1/z_j\in\mathbb{H}
\)。
Then the endpoint Poisson weight and pseudo-hyperbolic distance satisfy
\[
P_{a_j}(-1)=\Im(t_j),
\qquad
\rho(a_j,a_k)=\left|\frac{t_j-t_k}{t_j-\overline t_k}\right|.
\]
Define the upper half-plane Dirichlet Green function
\[
G_{\mathbb{H}}(t,s):=\log\left|\frac{t-\overline s}{t-s}\right|.
\]
Then generalized defect entropy (negative spectral volume potential)
\(
S_{\mathrm{gen}}:=-\log\det\mathsf P
\)
has strict equivalence
\[
\boxed{\;
S_{\mathrm{gen}}
=\sum_{j=1}^{\kappa}\bigl(-\log\Im(t_j)\bigr)
+2\sum_{1\le j<k\le\kappa}G_{\mathbb{H}}(t_j,t_k)\; }.
\]
And if the defective multiset is decomposed into a disjoint union \(A\sqcup B\), then the intersection can satisfy the identity
\[
\boxed{\;
-\log\det\mathsf P(A\cup B)
=-\log\det\mathsf P(A)-\log\det\mathsf P(B)
+2\sum_{a\in A}\sum_{b\in B}\bigl(-\log\rho(a,b)\bigr)\; }.
\]
This leads to the submultiplicative law \(\det\mathsf P(A\cup B)\le \det\mathsf P(A)\det\mathsf P(B)\), and the cross term is completely controlled by the pseudo-hyperbolic separation.
\end{conclusion}
\begin{proof}[Proof sketch]
For the image charge mirror energy representation, see proposition \ref{prop:xi-pick-poisson-image-charge-green-energy};
For the identity and submultiplicativity of union cross mutual energy, see the proposition \ref{prop:xi-pick-poisson-union-submultiplicativity-cross-energy}.
\end{proof}

\begin{conclusion}[Schur complementary chain externalization accounting: the new energy added point by point is equal to the conditional flux loss]\label{con:xi-terminal-pick-poisson-schur-ledger}
For a fixed point sequence \(1,\dots,\kappa\), let \(\mathsf P^{(m)}\) be the principal submatrix induced by the first \(m\) points, and define the Schur complement (conditional flux)
\[
\pi_m
:=\mathsf P_{mm}-\mathsf P_{m,1:m-1}\,(\mathsf P^{(m-1)})^{-1}\,\mathsf P_{1:m-1,m},
\qquad
m=1,\dots,\kappa.
\]
then there is a strict decomposition
\[
\boxed{\;
\det\mathsf P=\prod_{m=1}^{\kappa}\pi_m,
\qquad
S_{\mathrm{gen}}=-\sum_{m=1}^{\kappa}\log\pi_m\; }.
\]
And every \(\pi_m\) satisfies \(0<\pi_m\le \mathsf P_{mm}=P_{a_m}(-1)\).
Therefore, the negative spectral volume potential has a stepwise accounting structure: the new cost of the \(m\)th defect is exactly equal to its conditional flux loss on the spanned space of existing defects, and is controlled by an explicit upper bound on the endpoint self-energy.
\end{conclusion}
\begin{proof}[Proof sketch]
See proposition \ref{prop:xi-pick-poisson-schur-flux-ledger}.
\end{proof}

\begin{conclusion}[Spectral margin and pathological sealing:\texorpdfstring{$\lambda_{\min}(\mathsf P)$}{lambda_min(P)} collapses simultaneously with the condition number in the welt/collision limit]\label{con:xi-terminal-pick-poisson-conditioning-collapse}
For the Pick--Poisson matrix \(\mathsf P(t)\) induced by the scan spectrum point set (proposition \ref{prop:xi-scan-toeplitz-conditioning-by-separation}), record the endpoint weight \(p_j(t):=P_{a_j(t)}(-1)\) and the pseudo-hyperbolic distance \(\rho_{jk}(t):=\rho(a_j(t),a_k(t))\).
Then its minimum eigenvalue satisfies explicit clamping: on the one hand
\[
\boxed{\;
\lambda_{\min}(\mathsf P(t))
\ \ge\
\frac{\Bigl(\prod_{j=1}^{\kappa}p_j(t)\Bigr)\Bigl(\prod_{1\le j<k\le\kappa}\rho_{jk}(t)^2\Bigr)}
{\Bigl(\sum_{j=1}^{\kappa}p_j(t)\Bigr)^{\kappa-1}}\; } ;
\]
On the other hand, let \(p_{\max}(t):=\max_j p_j(t)\), \(\rho_{\min}(t):=\min_{j<k}\rho_{jk}(t)\), then there is a quadratic upper bound
\[
\boxed{\;
\lambda_{\min}(\mathsf P(t))\ \le\ p_{\max}(t)\,\rho_{\min}(t)^2\; }.
\]
Then from \(\lambda_{\max}(\mathsf P)\le \mathrm{tr}(\mathsf P)\) and inference the lower bound of \ref{cor:xi-pick-poisson-gap-lower-bound}, we can get the condition number seal
\[
\boxed{\;
\kappa(\mathsf P)
:=\frac{\lambda_{\max}(\mathsf P)}{\lambda_{\min}(\mathsf P)}
\ \le\
\frac{\bigl(\mathrm{tr}\,\mathsf P\bigr)^{\kappa}}
\Bigl(\prod_{j}P_{a_j}(-1)\Bigr)\Bigl(\prod_{j<k}\rho(a_j,a_k)^2\Bigr)\; }.
\]
Therefore when \(\rho_{\min}(t)\downarrow 0\) (collision) or \(p_j(t)\downarrow 0\) (edge), the spectral margin on the Pick/Toeplitz side must collapse at at least a quadratic rate, resulting in an auditable ill-posed certificate.
\end{conclusion}
\begin{proof}[Proof sketch]
See proposition \ref{prop:xi-scan-toeplitz-conditioning-by-separation} for upper and lower bounds;
For sealing of condition number fractions see Corollary \ref{cor:xi-pick-poisson-condition-number-certificate} (and use the closed form decomposition of \(\det\mathsf P\)).
\end{proof}

\begin{conclusion}[Hankel--Toeplitz common ill conditioning: the minimum stability constant shared by two types of inversions is Vandermonde's square law]\label{con:xi-terminal-hankel-toeplitz-common-illconditioning}
Under the finite atomic moment aperture \(\mu=\sum_{j=1}^{\kappa}w_j\delta_{a_j}\), note \(b_j:=a_j^{-1}\in(0,1)\), \(\widetilde w_j:=w_j b_j\), and let Hankel main block \(\mathsf H_{\kappa-1}:=(\mu_{p+q})_{0\le p,q\le \kappa-1}\).
Then there is a closed determinant
\[
\boxed{\;
\det \mathsf H_{\kappa-1}
=\left(\prod_{j=1}^{\kappa}\widetilde w_j\right)
\left(\prod_{1\le j<k\le \kappa}(b_k-b_j)^2\right)\; }.
\]
Therefore any error amplification factor based on Hankel inversion (Prony/Pad\'e/Hankel kernel space method) cannot be consistently bounded in the \(b\)-collision limit, at least as divergent as \(\prod_{j<k}|b_k-b_j|^{-2}\).
The Vandermonde square law is the same as
\(
\det\mathsf P=(\prod_j P_{a_j}(-1))(\prod_{j<k}\rho(a_j,a_k)^2)
\)
The degradation mechanism is of the same type, indicating that the "moment certificate channel (Toeplitz/Pick)" and the "tomographic inversion channel (Hankel/Vandermonde)" have homologous pathologies in the collision parameters, and a unified lower-bound stability guardrail can be given by the determinant closed form.
\end{conclusion}
\begin{proof}[Proof sketch]
Vandermonde decomposition and pathological divergence of Hankel determinant see proposition \ref{prop:xi-depth-hankel-determinant-vandermonde-square};
Pick--Poisson's product law of the same type is shown in the proposition \ref{prop:xi-pick-poisson-det-factorization}.
\end{proof}

\begin{conclusion}[Toeplitz--Minimum failure location of PSD certificate: equivalent to CMV reflection spectrum prefix out of bounds]\label{con:xi-terminal-toeplitz-psd-failure-cmv-prefix}
Let \(\mathbf{T}_m\) be the Toeplitz principal submatrix generated by the horizon Carath\'eodory moment \((c_n)\), and let \((\alpha_n)_{n\ge 0}\) be the reflection coefficient (Verblunsky parameter) given by the corresponding Schur algorithm.
Then the failure of the Toeplitz--PSD level can be accurately located by the minimum cross-boundary reflection coefficient: there is a minimum \(m^\ast\ge 1\) such that
\[
\boxed{\ |\alpha_{m^\ast-1}|\ge 1\ },
\]
And this \(m^\ast\) equivalently describes the "minimum PSD failure order" (still maintains \(\mathbf{T}_m\succeq 0\) at lower orders \(m<m^\ast\)).
The positions of falsifier occurrences are therefore condensed into auditable minimum prefix metrics without relying on global numerical spectral ordering.
\end{conclusion}
\begin{proof}[Proof sketch]
See the theorem \ref{thm:xi-toeplitz-det-verblunsky} (Toeplitz determinant reflection product solution with minimum failure index) and the discussion that follows.
\end{proof}

\begin{conclusion}[Chebyshev--integer coefficient of Frobenius \texorpdfstring{$\delta_p$}{delta_p}-structure: Adams double angle becomes Frobenius-lift]\label{con:xi-terminal-chebyshev-frobenius-pderivation}
At boundary coordinates \(S=2\cos(\theta/2)\), the power map \(w\mapsto w^{m}\) is equivalent to \(S\mapsto C_m(S)\) (Chebyshev--Adams polynomial).
Define difference polynomials for prime numbers \(p\)
\[
\Delta_p(S):=\frac{C_p(S)-S^p}{p}\in\ZZ[S],
\]
And define for any \(f\in\ZZ[S]\)
\[
\delta_p(f):=\frac{f(C_p(S))-\bigl(f(S)\bigr)^p}{p}.
\]
Then there are integral coefficient closure and Frobenius-lift identity
\[
\boxed{\ \delta_p(f)\in\ZZ[S],\qquad f(C_p(S))=\bigl(f(S)\bigr)^p+p\,\delta_p(f)\ },
\]
And the modulo \(p\) level strictly degenerates into \(f(C_p(S))\equiv f(S)^p\pmod p\).
Therefore the completion function of multiple angle/power semi-groups can be promoted as the algebraic source of \(p\)-advanced congruence guardrails.
\end{conclusion}
\begin{proof}[Proof sketch]
See definition \ref{def:chebyshev-frobenius-pderivation} and proposition \ref{prop:chebyshev-frobenius-pderivation-integrity}.
\end{proof}

\begin{conclusion}[Witt telescopic Koenigs linearization: \texorpdfstring{$s\mapsto ms$}{s->ms} strictly translated on the depth line \texorpdfstring{$y$}{y}]\label{con:xi-terminal-witt-koenigs-logm-translation}
Under the Dirichlet--Abel interface \(z=\lambda^{-s}\), Witt scaling \(s\mapsto ms\) is strictly equivalent to power substitution \(z\mapsto z^m\).
Define parameterless depth coordinates
\[
y(s):=\log\!\bigl(s\log\lambda\bigr)=\log\!\bigl(-\log(\lambda^{-s})\bigr),
\]
Then there is a strict identity for any integer \(m\ge 1\)
\[
\boxed{\ y(ms)=y(s)+\log m\ }.
\]
Therefore the Witt stretching is completely diagonalized into a translational semigroup in \(y\)-coordinates; the corresponding Witt operator on the depth line \(y\) becomes a translational superposition operator generated by the discrete measure \(\nu=\sum_{m\ge 1}\frac{1}{m}\delta_{\log m}\).
\end{conclusion}
\begin{proof}[Proof sketch]
See theorem \ref{thm:koenigs-linearization-witt-dilation} and corollary \ref{cor:witt-depth-line-convolution}.
\end{proof}

\begin{conclusion}[Local uplift selection law of window-\texorpdfstring{$6$}{6}: \texorpdfstring{$SO(10)$}{SO(10)} branch is the only top-level candidate that preserves locality]\label{con:xi-terminal-window6-local-uplift-so10-unique}
Under the strong local uplift caliber implemented by \(\mathrm{Aut}(Q_6)\) keeping the stable type label \(F_6\) constant, only \(\delta=\pm 34\) is allowed in addition to trivial displacements (corollary \ref{cor:terminal-window6-local-uplift-admissibility}).
On the other hand, the tail three-branch set of window-\(6\) uplift is \(\{21,34,55\}=\{F_8,F_9,F_{10}\}\), and is aligned with the top-level Zeckendorf terms of \((SU(5),SO(10),E_6)\) respectively as \((F_8,F_9,F_{10})\) (Theorem \ref{thm:terminal-window6-tail-three-branch}).
Therefore, when additionally imposing the protocol constraint that "uplift must be locally achievable and label unchanged", the \(\delta=34\) branch (\(F_9\)) is the only top-level completion that is compatible with locality; the top-level alignment of the other two branches can only be achieved externally through non-local mechanisms.
\end{conclusion}
\begin{proof}[Proof sketch]
For local realizability and uniqueness of \(\delta=\pm 34\) see corollary \ref{cor:terminal-window6-local-uplift-admissibility};
The three-branch alignment with the Zeckendorf--GUT top term is shown in theorem \ref{thm:terminal-window6-tail-three-branch} (and is consistent with the selection of \(F_9\) of the family number locking law \ref{thm:terminal-family-uplift-lock}).
\end{proof}

\begin{conclusion}[bin-fold last bit sufficiency of forward distribution: total variation collapse and single parameter exponential skew]\label{con:xi-terminal-foldbin-lastbit-sufficient-statistic}
Let \(\mu_m(w):=d_m^{\mathrm{bin}}(w)/2^m\) be the forward distribution of bin-fold on \(X_m\), and let the last bit indicate \(\ell(w):=w_m\in\{0,1\}\).
Divide \(X_m\) into blocks according to the last position and let \(\overline\mu_m\) be the homogenization approximation within the block, then there is an absolute constant \(C>0\) such that
\[
\boxed{\ D_{\mathrm{TV}}(\mu_m,\overline\mu_m)\ \le\ C\left(\frac{\varphi}{2}\right)^{m}\ }.
\]
And there exists \(\beta_m\in\RR\) such that \(\overline\mu_m\) can be written as the exponential slope of the stable layer uniform distribution \(u_m\)
\[
\overline\mu_m(w)=Z_m^{-1}u_m(w)e^{-\beta_m\ell(w)},
\qquad
\beta_m\to\log\varphi.
\]
Therefore, the last bit \(\ell(w)\) becomes a sufficient statistic with exponential precision, and the bin-fold output distribution is compressed into a single-parameter exponential family in the sense of information geometry.
\end{conclusion}
\begin{proof}[Proof sketch]
See proposition \ref{prop:fold-bin-lastbit-sufficient-tv}.
\end{proof}

\begin{conclusion}[Two-state closure of escort temperature family and divergence spectrum: Bernoulli degradation and Rényi entropy rate collapse of KL between temperatures]\label{con:xi-terminal-foldbin-escort-divergence-two-state-closure}
Define the escort distribution \(\pi_{m,q}(w)\propto (d_m^{\mathrm{bin}}(w))^{q}\) for \(q\ge 0\), and let
\(
\theta(q):=\lim_{m\to\infty}\mathbb{P}_{\pi_{m,q}}(w_m=1)
=\frac{1}{1+\varphi^{q+1}}
\)。
Then for any \(p,q\ge 0\), the inter-temperature KL limit exists and strictly degenerates into one-dimensional Bernoulli KL:
\[
\boxed{\;
\lim_{m\to\infty}D_{\mathrm{KL}}(\pi_{m,q}\Vert\pi_{m,p})
=\theta(q)\log\!\frac{\theta(q)}{\theta(p)}+(1-\theta(q))\log\!\frac{1-\theta(q)}{1-\theta(p)}\; }.
\]
At the same time, the Rényi divergence limit of the bin-fold forward distribution \(\mu_m\) relative to the uniform baseline of the stable layer exists and is given in a closed form for every \(\alpha>0\), with exponential convergence of \(O((\varphi/2)^m)\) (Theorem \ref{thm:fold-bin-renyi-divergence-limit}).
And the Rényi entropy density satisfies uniform collapse: for all \(q>0\)
\[
\boxed{\ \lim_{m\to\infty}\frac{1}{m}H_q(\mu_m)=\log\varphi\ },
\]
Thus the "multi-temperature/multi-divergence/multi-fractal" structure is institutionally compressed in the limit to a one-dimensional spectrum of the last Bernoulli parameter \(\theta(q)\).
\end{conclusion}
\begin{proof}[Proof sketch]
Bernoulli KL degeneration see proposition \ref{prop:fold-bin-escort-escort-kl};
See the theorem for Rényi divergence limit \ref{thm:fold-bin-renyi-divergence-limit};
Rényi entropy rate collapse theorem \ref{thm:fold-bin-renyi-rate-collapse}.
\end{proof}

\begin{conclusion}[bin-fold Linear-time computability of fiber multiplicity: Fibonacci-digit DP gives a recomputable audit kernel]\label{con:xi-terminal-foldbin-digit-dp-linear-time}
Given \(m\ge 1\) and \(w\in \(Z(\mathrm{slack})=(b_1,\dots,b_K)\).
Then \(d_m^{\mathrm{bin}}(w)\) can be determined by the state
\(
\mathrm{DP}(k,\mathrm{prev},\mathrm{tight})
\)
The digit DP is calculated accurately in \(O(K-m)=O(m)\) time; only limited conflict constraints (no adjacent \(11\)), tightness constraints and last digit constraints \(w_m=1\Rightarrow c_{m+1}=0\) are imposed in the transfer.
Therefore, the fiber multiplicity of bin-fold not only has an auditable asymptotic law, but also has an accurate calculation kernel of linear complexity, providing algorithmic closure for fully deterministic recalculation of large windows.
\end{conclusion}
\begin{proof}[Proof sketch]
See proposition \ref{prop:fold-bin-digit-dp}.
\end{proof}

\endinput



